\documentclass[11pt]{amsart}
\usepackage{amsmath,amssymb,mathtools}
\usepackage[margin=1.18in]{geometry}
\usepackage{color}
\definecolor{paperblue}{rgb}{0.06,0.20,0.42}
\usepackage[bookmarksdepth=2,colorlinks=true,linkcolor=paperblue,citecolor=paperblue,urlcolor=paperblue]{hyperref}
\hypersetup{pdftitle={Beilinson-Deligne Periods and Gross-Stark Regulators of 1-Motives},
            pdfauthor={Mohammadreza Mohajer}}
\raggedbottom
\numberwithin{equation}{section}
\newtheorem{theorem}{Theorem}[section]
\newtheorem{proposition}[theorem]{Proposition}
\newtheorem{corollary}[theorem]{Corollary}
\newtheorem{lemma}[theorem]{Lemma}
\theoremstyle{definition}
\newtheorem{definition}[theorem]{Definition}
\newtheorem{example}[theorem]{Example}
\theoremstyle{remark}
\newtheorem{remark}[theorem]{Remark}
\newcommand{\Q}{\mathbf Q}
\newcommand{\Z}{\mathbf Z}
\newcommand{\R}{\mathbf R}
\newcommand{\C}{\mathbf C}
\newcommand{\Qbar}{\overline{\mathbf Q}}
\newcommand{\Qp}{\mathbf Q_p}
\newcommand{\Cp}{\mathbf C_p}
\newcommand{\Gm}{\mathbf G_m}
\newcommand{\Ga}{\mathbf G_a}
\newcommand{\OO}{\mathcal O}
\newcommand{\Res}{\operatorname{Res}}
\newcommand{\Gal}{\operatorname{Gal}}
\newcommand{\Tr}{\operatorname{Tr}}
\newcommand{\ord}{\operatorname{ord}}
\newcommand{\Fr}{\operatorname{Fr}}
\newcommand{\Lie}{\operatorname{Lie}}
\newcommand{\Hom}{\operatorname{Hom}}
\newcommand{\End}{\operatorname{End}}
\newcommand{\im}{\operatorname{im}}
\newcommand{\coker}{\operatorname{coker}}
\newcommand{\rank}{\operatorname{rank}}
\newcommand{\gr}{\operatorname{gr}}
\newcommand{\Irr}{\operatorname{Irr}}

\title[Special values and regulator relations of 1-motives]
{Beilinson--Deligne Periods and Gross--Stark Regulators of 1-Motives}
\author{Mohammadreza Mohajer}
\subjclass[2020]{Primary 11R42, 11R23; Secondary 14C15, 11G40, 11J86, 14F30}
\keywords{1-motive, Deligne period, mixed Euler factor, relative norm torus,
Gross--Stark regulator, local reciprocity, $p$-adic $L$-function, Artin induction}

\begin{document}

\begin{abstract}
For a Deligne 1-motive, the valuations of the lattice points
determine the difference between its local Euler factors and those
of its weight-graded pieces. We use this calculation to prove
Scholl's critical-value conjecture for every critical constant
lattice in a product of relative norm tori and split tori. We give
the exact order of vanishing and a relative $p$-adic residue formula.
For the odd $p$-units of a Galois CM field, local Tate duality
identifies the normalized Gross matrix with the obstruction to
lifting local characters to a Cartier-dual realization. Its Artin
blocks determine the linear relations among regulator entries and
the normalized Gross--Stark coefficients of order one. For the
full Deligne--Ribet functions, we prove that all polynomial
relations over the field generated by the cyclotomic variable and
rational Euler factors come from equalities of induced Artin
characters. Specialization at zero retains only constituents with
decomposition-group invariants. In an explicit $S_3$ example at
$p=7$, a constituent absent from the leading-coefficient identity gives
a quotient of normalized germs with derivative $42$ modulo $49$.
\end{abstract}

\maketitle
\setcounter{tocdepth}{1}
\tableofcontents
\clearpage

\section{Introduction}

A prime can enter the $L$-function of a 1-motive through the
valuation of a single rational point. Scholl's example
$[\Z\xrightarrow{1\mapsto q}\Gm]$ gives the simplest instance.
Its two graded pieces contribute $\zeta(s)\zeta(s+1)$. The
extension removes one inertia invariant at $q$, so the mixed
$L$-function acquires the factor $1-q^{-s}$. The pole at $s=0$
then cancels, leaving $-\tfrac12\log q$. The same $\log q$ is
the real period obtained by integrating $dx/x$ from $1$ to $q$
\cite[\S II]{Scholl}. The example raises two questions. Which
valuation spaces enter the formula for a general 1-motive? What
arithmetic information survives when the real logarithm is replaced
by a $p$-adic logarithm?

Let $M=[Y\to G]$ be a Deligne 1-motive over $\Q$, with $Y$ an
\'etale lattice and $G$ an extension of an abelian variety $A$ by
a torus $T$. The lattice, the abelian variety, and the torus give
the homological weights $0,-1,-2$. Write $T_B(M)$ and
$T_{\rm dR}(M)$ for its rational Betti and de Rham realizations.
The comparison at infinity induces
\begin{equation}\label{eq:intro-comparison}
 \alpha_M:T_B(M)^+\otimes_\Q\R\longrightarrow
                  (T_{\rm dR}(M)/F^0)\otimes_\Q\R .
\end{equation}
The plus sign denotes invariants under complex conjugation.
Scholl calls $M$ critical when $\alpha_M$ is an isomorphism and
writes $c^+(M)$ for its determinant in rational bases, up to
$\Q^\times$. His Conjecture~A predicts $L(M,0)/c^+(M)\in\Q$
\cite[\S II, Conj.~A]{Scholl}. Here the local factors use inertia
invariants of the full mixed realization. This is the sense in which
we use the Beilinson--Deligne critical-value conjecture. Deligne's
pure critical-value conjecture and Beilinson's conjectures on
motivic regulators supply its background
\cite{DeligneL,Beilinson,Nekovar}. When the mixed comparison has a
kernel, we specify the additional determinant map whose leading
term we compute.

\subsection*{From valuations to critical values}
Theorem~\ref{thm:euler} gives the missing Euler polynomial at every
finite prime $q$ in terms of a rational map
$\nu_q:Y_\Q\to X_*(T)_\Q$, including the case of ramified lattice
or torus. Proposition~\ref{prop:realdefect} describes the
archimedean defect using the noncompact real directions of the
torus. For the norm-one torus of an extension $K/k$ and its lattice
of norm-one $S$-units, the finite directions are the relative
divisors above $S$. Appending their valuations weighted by $\log q$
to \eqref{eq:intro-comparison} gives a square comparison. Its
determinant is the leading coefficient up to $\Q^\times$
(Theorem~\ref{thm:relative}).

Theorem~\ref{thm:allcritical} goes further. Fix a product of
relative norm tori and a split torus of rank $t$, and allow
\emph{any} constant lattice map for which $M$ is critical.
The order of $L(M,s)$ at zero is $\sum_q\rank\nu_q-t$.
Criticality forces the sum of the ranks to be at least $t$.
When equality holds, the lattice vectors with zero valuation span
the ordinary relative units and the remaining vectors account for
all split real logarithms. This proves Scholl's Conjecture~A for
the entire family. For totally real fields the corresponding
$p$-adic determinant comes from Colmez's residue formula
(Theorem~\ref{thm:prelative} and Corollary~\ref{cor:splitpadic}).
Over a real quadratic field, Theorem~\ref{thm:rigidity} shows
that the ratio of the graded special value to the period is algebraic
at either place exactly when the point is a global unit.

\subsection*{From local reciprocity to Gross--Stark coefficients}
Let $N/\Q$ be a Galois CM field, a totally imaginary quadratic
extension of a totally real field. On the odd part of its
$p$-units, the valuations of local norms identify the unit space
with a divisor space. Transporting negative norm logarithms by this
identification gives an operator $R_N$. If $F\subset N$ is totally
real and $\chi$ is a totally odd character over $F$, restricting
$R_N$ to its $\chi^{-1}$-space gives the normalized Gross matrix
$B(F,\chi)$ (Proposition~\ref{prop:descent}). A residue-degree
factor for each split prime ensures compatibility across base fields.

By Theorem~\ref{thm:cuplift}, local Tate duality recovers the
matrix itself. A linear combination of local valuation and logarithm
characters lifts to the dual realization precisely when its two
coefficient vectors satisfy $x=-B(F,\chi)^{\mathsf t}y$.
It extends the use of the adjoint Kummer map in
\cite[eq.~(2.3), Prop.~7.5]{MohajerSemistable} to the simultaneous
character projections needed here. Varying the local directions
recovers every principal minor. At the cyclotomic direction,
the determinant is a normalized Taylor coefficient $A(F,\chi)$
by the theorem of Dasgupta, Kakde and Ventullo \cite{DKV}.

The common operator identifies the linear relations among entries
from different fields (Theorem~\ref{thm:joint}). For coefficients
of expected order one, it gives their affine relations and exact
local fields of definition
(Corollaries~\ref{cor:rankone} and~\ref{cor:coefficientfield}).
For arbitrary orders it gives unconditional multiplicative
identities. Weak $p$-adic Schanuel identifies the complete ideal of
relations among the numerical coefficients
(Theorem~\ref{thm:ideal}). The independence of logarithms for a
single irreducible regulator matrix and quantitative rank bounds
are present in Maksoud's work \cite[Thms.~1.5, 4.3]{Maksoud}.
Our comparison keeps track of the same irreducible block when it
occurs in several fields. Hofer and Kleine prove exact orders in
families of unbounded rank \cite{HK}. In our examples the individual
coefficients in such families are expressed in common rank-one
factors.

\subsection*{What the full \texorpdfstring{$p$}{p}-adic functions remember}
Write $\mathcal H_{F,\chi}(s)$ for the Deligne--Ribet function
after division by its expected power of $s$ and the explicit
constant in \eqref{eq:normcoeff}. Its value at zero is $A(F,\chi)$.
Theorem~\ref{thm:germrelations} determines \emph{all} polynomial
relations among the analytic germs for a finite family. Over the
field generated by $s$, $p$-adic constants and the cyclotomic powers
of positive rational $p$-units, the relations are generated by
equalities of direct sums of the full induced Artin representations.
The proof compares the interpolation values at integers which tend to
zero $p$-adically and to minus infinity in the real absolute value.
The complex functional equation and Chebotarev's theorem then
recover the full characters. Brauer induction gives algebraic
independence of the meromorphic $p$-adic functions of distinct
irreducible totally even Artin representations over $\Q$
(Corollary~\ref{cor:artingerms}). No transcendence conjecture enters
this argument.

Only the constituents with a fixed vector under the decomposition
group at $p$ enter the normalized constant. In the explicit
nonabelian example of Section~\ref{sec:examples}, a quadratic
constituent with no such vector has normalized value one. Its
normalized first derivative is $42$ modulo $49$, giving the precise first-order
failure of a leading-coefficient identity
(Proposition~\ref{prop:next}).

The additive period theorems for 1-motives treat a related question.
Huber and W\"ustholz prove the complex theorem on linear relations
of 1-periods \cite{HW}. Mohajer and Sebbar study comparison spaces and
linear relations in good reduction \cite{MS,MohajerThesis}. Here the
semistable Kummer logarithms furnish matrix entries, and special
values involve their determinants. Section~\ref{sec:limits} explains
where abelian periods, elliptic exceptional zeros and height
regulators enter for 1-motives with an abelian part.

\section{Local Euler factors and real criticality}
\label{sec:local}

We work up to isogeny, with rational homological realizations.
The Tate twist $\Q_\ell(1)$ has eigenvalue $q^{-1}$ under
geometric Frobenius at $q\ne\ell$. For a number field $L$ and
prime $\mathfrak w$ of $L$, $\ord_{\mathfrak w}$ denotes the
discrete valuation taking value $1$ on a uniformizer. A
subscript $S$ on units or divisors refers to all primes above
the rational primes of $S$. Tensor products of multiplicative
groups with characteristic-zero fields are written additively.

\subsection{The valuation of a point on a semiabelian variety}

Let $F$ be a finite extension of $\Q_q$ and let
$0\to T\to G\xrightarrow{\pi} A\to0$ be a semiabelian variety.
First assume $T$ splits over $F$. Pairing with characters
defines $v_{T,F}:T(F)\to X_*(T)_\Q$, normalized using the
valuation $v_F=\ord_F/\ord_F(q)$. There is a unique homomorphism
\begin{equation}\label{eq:semival}
 v_{G,F}:G(F)\longrightarrow X_*(T)_\Q
\end{equation}
extending $v_{T,F}$ and killing compact open subgroups.
In fact choose a compact open subgroup $C$ of $G(F)$ small
enough to be pro-$q$. The image $\pi(C)$ is open in $A(F)$
by smoothness and the local submersion theorem, hence has
finite index since $A(F)$ is compact. For $g\in G(F)$, take
$n>0$ and $c\in C$ with $\pi(ng)=\pi(c)$. Set
$v_{G,F}(g)=n^{-1}v_{T,F}(ng-c)$.
The valuation on $T(F)$ kills $T(F)\cap C$ because the image
of this compact subgroup under the discrete valuation is
finite and torsion free. This gives independence of $c$.
Passing to a common multiple gives independence of $n$,
and the same argument proves additivity. Two compact open
subgroups have an open intersection of finite index in each,
so the map kills every compact subgroup. The defining
properties give uniqueness. Normalizing by $v_F(q)=1$
makes the map compatible with extension to any finite field
over which the torus is split. Indeed every compact subgroup
remains compact after extension, and the valuation on the
torus restricts to the original valuation.

Choose a finite Galois extension $F/\Q_q$ on which $Y$ and $T$
split. The points $u(Y)$ then belong to $G(F)$. Define
$\nu_q:Y_\Q\to X_*(T)_\Q$ by $v_{G,F}\circ u$ and
scalar extension. Uniqueness of \eqref{eq:semival} implies
Galois equivariance. Compatibility under further finite
extension makes $\nu_q$ independent of $F$ once a geometric
local base is fixed. These maps vanish for all but finitely
many $q$. For this last assertion, spread out the split
semiabelian variety and the finitely many images of lattice
generators over the integers of a number field with finitely
many primes inverted. At each remaining place those images
are integral points of a smooth group model of finite type.
The integral points form a compact subgroup.

\subsection{The local correction}

Write $V_\ell M=T_\ell M\otimes\Q_\ell$, and similarly for
$Y,T,A$. For $q\ne\ell$, let $P_q(B,t)$ denote
\[
 P_q(B,t)=\det(1-t\Fr_q\mid(V_\ell B)^{I_q})
\]
where $I_q$ is inertia at $q$. This notation for $Y$ uses
the Artin representation $Y_{\Q_\ell}$.

\begin{theorem}\label{thm:euler}
For any 1-motive $M=[Y\xrightarrow{u}G]$ over $\Q$ there is an exact
sequence of geometric Frobenius modules
\begin{equation}\label{eq:euler-seq}
 0\longrightarrow (V_\ell G)^{I_q}
 \longrightarrow (V_\ell M)^{I_q}
 \longrightarrow (\ker\nu_q)^{I_q}\otimes_\Q\Q_\ell
 \longrightarrow0 .
\end{equation}
Consequently, with $D_q(M)=(\im\nu_q)^{I_q}$,
\begin{equation}\label{eq:localeuler}
 P_q(M,t)
 =
 \frac{P_q(Y,t)P_q(A,t)P_q(T,t)}
 {\det(1-t\Fr_q\mid D_q(M))}.
\end{equation}
In particular the correction polynomial is rational and independent
of $\ell$. For constant $Y$, it is $(1-t)^{\rank\nu_q}$.
\end{theorem}

\begin{proof}
For any abelian variety $B/F$, the rational $\ell$-adic Kummer
class of each point of $B(F)$ is zero. Take a compact open pro-$q$
subgroup on which multiplication by $\ell$ is an automorphism.
Every point has a positive multiple in this subgroup, where
compatible $\ell$-power divisions exist over $F$. Dividing the
resulting identity in $H^1(F,V_\ell B)$ proves the assertion.
The same argument shows that the points of a small pro-$q$
subgroup of $G(F)$ have zero rational Kummer class.

We first split $0\to V_\ell T\to V_\ell G\to V_\ell A\to0$
as a representation of the absolute Galois group of $F$.
Deligne's Cartier duality identifies the dual of $[0\to G]$
with $[X^*(T)\to A^\vee]$, with Tate realizations paired
into $\Q_\ell(1)$ \cite[\S10.2]{DeligneIII}. Over $F$ the
lattice $X^*(T)$ is constant. The extension classes of
the Tate realization of the dual are the Kummer classes
of its lattice images in $A^\vee(F)$, which vanish by the
preceding argument. Dualizing gives the desired splitting.
Averaging a splitting over $\Gal(F/\Q_q)$ gives a splitting
over $\Q_q$ because coefficients have characteristic zero.
In particular $P_q(G,t)=P_q(T,t)P_q(A,t)$, and the map from
$H^1(I_F,V_\ell T)$ to $H^1(I_F,V_\ell G)$ is injective.

The connecting map for $0\to V_\ell G\to V_\ell M
\to Y_{\Q_\ell}\to0$, restricted to $I_F$, sends each
$y$ to the restricted Kummer class of $u(y)$.
Write $nu(y)=t+c$ with $t\in T(F)$ and $c$ in the chosen
pro-$q$ subgroup. The Kummer class of $c$ is zero, and after
restriction the class of $t$ is given by its vector of
valuations times the $\ell$-adic tame character.
For $\Gm$ this follows by taking $\ell$-power roots of the unit
part over the maximal unramified extension, where the
residue field is algebraically closed. A uniformizer gives
a nonzero Kummer character on inertia because its $\ell^n$-th root
has valuation of denominator $\ell^n$ in the integral valuation
normalization. The denominators are unbounded, so the resulting
$\ell$-adic tame character is nonzero even after rationalization. Applying the split torus characters proves the claim
for $T$. Injectivity for $T\to G$ now identifies the kernel
of the connecting map over $I_F$ with
$\ker\nu_q\otimes\Q_\ell$. The map
$H^1(I_q,V_\ell G)\to H^1(I_F,V_\ell G)$ is injective by
inflation--restriction, as the quotient of inertia groups
is finite and the coefficients have characteristic zero.
Taking the kernel on $(Y_{\Q_\ell})^{I_q}$ yields
\eqref{eq:euler-seq}. Inertia acts through a finite group on
$Y_\Q$ and $X_*(T)_\Q$, so invariants commute with the image
of $\nu_q$. Taking determinants gives \eqref{eq:localeuler}.
\end{proof}

Let $L_{\mathrm{inv}}(M,s)$ denote the Euler product using
$P_q(M,q^{-s})^{-1}$ and let $L_{\gr}(M,s)$ denote the product
for the three weight pieces. The standard local factors of
abelian varieties, tori and Artin representations in this
formula are independent of $\ell$. For the abelian part,
the rational inertia invariants specialize to the Tate
module of the special fiber of the N\'eron model. The
formal kernel of specialization is uniquely $\ell$-divisible
and the component group is finite. The unipotent part
of the connected special fiber contributes no rational
$\ell$-adic Tate module. The torus and abelian quotient
over the residue field have compatible local polynomials.
For an abelian variety the homological convention shifts the usual
cohomological $L$-function by one, using a polarization.
Theorem~\ref{thm:euler} gives
\begin{equation}\label{eq:globaleuler}
 L_{\mathrm{inv}}(M,s)=L_{\gr}(M,s)
 \prod_q\det(1-q^{-s}\Fr_q\mid D_q(M)).
\end{equation}
The correction is a finite product. At zero its contribution
to the order of vanishing is
$\sum_q \dim_\Q D_q(M)^{\Fr_q=1}$, with a contribution
$(\log q)^{\dim D_q(M)^{\Fr_q=1}}$ to the leading coefficient
and, on the complementary Frobenius summand, the determinant
of $1-\Fr_q$. This follows from the finite order of $\Fr_q$
on $D_q(M)$ and the first-order expansion of $1-q^{-s}$.

\begin{example}
Let $J/\Q$ be quadratic, let $\eta$ be its sign character,
and twist the split Kummer 1-motive
$[\Z\xrightarrow{1\mapsto q}\Gm]$ by $\eta$ on both terms.
The twist descends because negation of the lattice and
inversion on the torus commute with the Kummer map.
Over a splitting field the valuation at $q$ is nonzero,
so its image spans the sign line in the cocharacters.
If $q$ is inert in $J$ then the correction polynomial
in \eqref{eq:localeuler} is $1+t$. In
\eqref{eq:globaleuler} it contributes the factor
$1+q^{-s}$, which has value $2$ at zero. If $q$
ramifies in $J$, inertia has zero invariants on the
sign line and the correction polynomial is $1$.
\end{example}

\subsection{The defect at the real place}

Put $\mathfrak g_\R=\Lie(G)\otimes_\Q\R$ and use
$\mathfrak t_\R,\mathfrak a_\R$ for $T,A$. Deligne's
realization via the universal vector extension identifies
$T_{\mathrm{dR}}(M)/F^0$ with $\Lie(G)$
\cite[\S10.1]{DeligneIII}. In the analytic description,
an integral Betti vector is a pair $(z,y)$ with
$\exp_G(z)=u(y)$. On this description the comparison
to the Lie algebra sends $(z,y)$ to $z$.
For the torus, let
\[
  C_T=\mathfrak t_\R/\alpha_T(T_B(T)^+\otimes_\Q\R).
\]
It is the real vector space of noncompact torus directions.

\begin{proposition}\label{prop:realdefect}
The natural map $C_T\to
\mathfrak g_\R/\alpha_G(T_B(G)^+\otimes\R)$ is an isomorphism.
There is a canonical exact sequence
\begin{equation}\label{eq:archdefect}
 0\to\ker\alpha_M\to Y_\Q^+\otimes\R
 \xrightarrow{\lambda_\infty} C_T
 \to\coker\alpha_M\to0 .
\end{equation}
The middle arrow is obtained by lifting a lattice vector to
the real Betti realization and taking its image modulo the
periods of $G$. In a product of a torus with an abelian variety,
it is the real logarithm of the toric lattice map, projected
to $C_T$.
\end{proposition}

\begin{proof}
The exact sequence of Betti realizations for $T,G,A$ stays
exact after taking plus parts and extending scalars to $\R$.
The corresponding Lie sequence is exact. The map for $T$ is
injective by the description $\ker\exp_T=2\pi iX_*(T)$.
For $A$ the map is an isomorphism because the period lattice
spans $\Lie(A)(\C)$ over $\R$ and comparison is compatible
with conjugation. Explicitly the real-linear logarithmic period map
$H_1(A(\C),\R)\to\Lie(A)(\C)$ is an isomorphism, so its restriction
to the invariant real subspaces is an isomorphism as well. The kernel-cokernel sequence of these
two exact rows gives injectivity of $\alpha_G$ and the
isomorphism of the asserted quotients. Applying the same
argument to $0\to T_B(G)^+\to T_B(M)^+\to Y_\Q^+\to0$
gives \eqref{eq:archdefect}. For $\Gm$ the plus projection
of a logarithmic lift is $\log|u(y)|$. The character
description of a torus gives the final assertion.
\end{proof}

\section{Relative norm tori and the complex leading term}
\label{sec:complex}

Let $K/k$ be a finite extension of number fields and let $S$
be a finite set of rational primes. Write
\begin{align*}
 T_{K/k}
 &=\ker\bigl(N_{K/k}:\Res_{K/\Q}\Gm\to\Res_{k/\Q}\Gm\bigr),\\
 U_{K/k,S}
 &=\ker\bigl(N_{K/k}:
 \OO_K[S^{-1}]^\times\otimes_\Z\Q\to
 \OO_k[S^{-1}]^\times\otimes_\Z\Q\bigr).
\end{align*}
The norm on split tori is a product along the fibers of the
map on embeddings. Its kernel is a torus. The inclusion
of the relative units gives a 1-motive up to isogeny,
\[
 M_S(K/k)=[U_{K/k,S}\longrightarrow T_{K/k}].
\]
In this notation the source means a constant lattice realizing
the indicated rational space. A free finite-index subgroup
of the integral norm-one $S$-units supplies such a lattice.
All determinant classes below are independent of the chosen
finite-index subgroup.

For each $q$ let $D_q(K/k)$ be the kernel of the map on
rational divisor spaces
\begin{equation}\label{eq:divnorm}
 \Q^{\{\mathfrak w\mid q\text{ in }K\}}
 \longrightarrow \Q^{\{\mathfrak v\mid q\text{ in }k\}},
 \qquad [\mathfrak w]\longmapsto
 f(\mathfrak w/\mathfrak v)[\mathfrak v].
\end{equation}
Set $d_q=\dim_\Q D_q(K/k)$ and $d_S=\sum_{q\in S}d_q$.
Taking valuations gives an exact sequence
\begin{equation}\label{eq:sunitseq}
 0\longrightarrow U_{K/k,\varnothing}
 \longrightarrow U_{K/k,S}
 \xrightarrow{(\ord_q)_{q\in S}}
 \bigoplus_{q\in S}D_q(K/k)\longrightarrow0 .
\end{equation}
Indeed, finiteness of the ideal class group gives
surjectivity of the valuation map for the absolute rational
$S$-units. For a vector in the relative divisor kernel,
lift it to an absolute $S$-unit of $K$. Its norm is an
ordinary unit of $k$ in the rationalized unit group.
Subtracting the image of that unit divided by $[K:k]$
gives a relative lift. The kernel consists of the
relative ordinary units. In particular
$b_S=\dim U_{K/k,S}=b_0+d_S$, where
\[
 b_0=r_1(K)+r_2(K)-r_1(k)-r_2(k).
\]
As usual, $r_1(L)$ and $r_2(L)$ count the real places and the
pairs of complex embeddings of $L$.

Define the real map
\begin{equation}\label{eq:augreal}
 \begin{split}
 \alpha_{M_S,S}:T_B(M_S)^+\otimes\R
 &\longrightarrow \mathfrak t_{K/k,\R}
       \oplus\bigoplus_{q\in S}\bigl(D_q(K/k)\otimes\R\bigr),\\
 x&\longmapsto \bigl(\alpha_{M_S}(x),
       ((\log q)\ord_q(\overline{x}))_{q\in S}\bigr).
 \end{split}
\end{equation}
Here $\overline{x}$ is the lattice image of $x$. We write
$c_S^+(M_S)$ for its determinant class when the map is an
isomorphism. Let $a_L=\operatorname{Res}_{s=1}\zeta_L(s)$
for the complex Dedekind zeta function. For a meromorphic
function, its leading term at zero is the coefficient of its
lowest nonzero Laurent power.

\begin{theorem}\label{thm:relative}
With this notation, the following statements hold.
\begin{enumerate}
\item The map \eqref{eq:augreal} is an isomorphism, and
\begin{equation}\label{eq:archformula}
 c_S^+(M_S)=
 \left[\frac{a_K}{a_k}\prod_{q\in S}(\log q)^{d_q}\right]
 \quad\text{in }\R^\times/\Q^\times .
\end{equation}
\item The two Euler products and the leading term satisfy
\begin{align}
 L_{\gr}(M_S,s)&=\zeta(s)^{b_S}
                 \frac{\zeta_K(s+1)}{\zeta_k(s+1)},\label{eq:grrelative}\\
 L_{\mathrm{inv}}(M_S,s)&=L_{\gr}(M_S,s)
                      \prod_{q\in S}(1-q^{-s})^{d_q},\label{eq:invrelative}\\
 \ord_{s=0}L_{\mathrm{inv}}(M_S,s)&=d_S,\nonumber\\
 L_{\mathrm{inv}}^*(M_S,0)&=
 \left(-\frac12\right)^{b_S}\frac{a_K}{a_k}
                 \prod_{q\in S}(\log q)^{d_q}.\label{eq:invrelativelead}
\end{align}
\end{enumerate}
In particular $M_\varnothing(K/k)$ is critical and
$[L(M_\varnothing,0)]=c^+(M_\varnothing)$.
\end{theorem}

\begin{proof}
We first compute the determinant for $S=\varnothing$.
For a number field $L$ the absolute unit 1-motive uses
$U_L=\OO_L^\times\otimes_\Z\Q$ and $\Res_{L/\Q}\Gm$.
Denote this motive by $M_L$. Augment its real comparison
by one rational vector $1/[L:\Q]\in L=\Lie(\Res_{L/\Q}\Gm)$.
\begin{equation}\label{eq:absaug}
 \bigl(T_B(M_L)^+\oplus\Q\bigr)\otimes\R
 \longrightarrow L\otimes_\Q\R,\qquad
 (x,t)\longmapsto \alpha_{M_L}(x)+t/[L:\Q].
\end{equation}
In coordinates $L\otimes\R\simeq\R^{r_1(L)}\times\C^{r_2(L)}$,
the plus Betti periods of the torus give one $2\pi i$
direction at each complex place. After projecting away
from those directions, a unit maps to its usual logarithmic
vector. We use $\log|\sigma(u)|$ at a real place and
$2\log|\sigma(u)|$ at a complex place. The unit vectors
span the sum-zero hyperplane by the Dirichlet unit theorem.
The last column of \eqref{eq:absaug} has trace one. This
proves invertibility.

Choose a $\Z$-basis of $\OO_L$ in the target and a basis of
the free part of the unit group in the source, with rational
Betti lifts. In the logarithmic coordinates the determinant
of the unit columns and the trace-one column is, up to sign,
the usual regulator $R_L$. This follows by replacing any
one coordinate by the sum of all coordinates and expanding
along that row. The factors of two at complex places in
the logarithmic coordinates affect the determinant only
by a nonzero rational number. The imaginary columns
contribute $(2\pi)^{r_2(L)}$. The change from an integral
basis of $\OO_L$ to real embedding coordinates has determinant
of absolute value $2^{-r_2(L)}\sqrt{|\operatorname{disc}L|}$.
The power of two comes from combining complex conjugate rows
in the discriminant matrix into their real and imaginary parts.
Consequently the
class of \eqref{eq:absaug} is
\[
 \left[(2\pi)^{r_2(L)}R_L/
                      \sqrt{|\operatorname{disc}L|}\right]=[a_L].
\]
The last equality is the analytic class number formula
\cite[Ch.~VII, \S5]{Neukirch}. Empty regulators are $1$.

The norm induces an exact sequence of rational 1-motives
$0\to M_\varnothing(K/k)\to M_K\to M_k\to0$.
On ordinary units, surjectivity after tensoring with
$\Q$ follows from $N_{K/k}(x)=x^{[K:k]}$ for $x\in k^\times$.
On the augmented real comparisons the norm is the field
trace in the target and is the identity on the auxiliary
$\Q$ in the source, since
$\Tr_{K/k}(1/[K:\Q])=1/[k:\Q]$. The kernel comparison is
$\alpha_{M_\varnothing(K/k)}$. Taking determinants in
these exact sequences gives the class $[a_K/a_k]$.

The inclusion of $M_\varnothing$ into $M_S$ induces an
exact sequence of plus Betti spaces with quotient
$\bigoplus_{q\in S}D_q(K/k)$. The quotient map is
\eqref{eq:sunitseq}. In bases respecting these subspaces,
\eqref{eq:augreal} has diagonal blocks
$\alpha_{M_\varnothing}$ and
$\bigoplus_{q\in S}(\log q)\mathrm{id}_{D_q(K/k)}$.
This proves \eqref{eq:archformula}.

The cocharacter module of $T_{K/k}$ is the kernel of the
map of two permutation modules on embeddings. At $q$, the
Frobenius cycles on the inertia orbits of the embeddings
have lengths $f(\mathfrak w/q)$ for $K$ and
$f(\mathfrak v/q)$ for $k$. Including the Tate twist thus
gives the zeta quotient in \eqref{eq:grrelative}.
The constant lattice contributes $\zeta(s)^{b_S}$.
In a local splitting field, the coordinate of the toric
valuation on an embedding inducing $\mathfrak w$ is
$\ord_{\mathfrak w}(u)/e(\mathfrak w/q)$, where
$e(\mathfrak w/q)$ is the ramification index. This embeds
the relative divisor space in the cocharacter space.
On each fiber over $\mathfrak v$ the sum of the embedding
coordinates is the $\mathfrak v$ coordinate of
\eqref{eq:divnorm} divided by $e(\mathfrak v/q)$.
By \eqref{eq:sunitseq}, the rank of $\nu_q$ is $d_q$
if $q\in S$, and is zero otherwise. Theorem~\ref{thm:euler}
proves \eqref{eq:invrelative}. The two simple zeta poles at
$1$ cancel. Using $\zeta(0)=-1/2$ and
$1-q^{-s}=s\log q+O(s^2)$ proves \eqref{eq:invrelativelead}.
For $S=\varnothing$ the lattice map has zero Kummer
boundary on every inertia group in characteristic
different from $\ell$. The weight filtration is split on
those inertia groups, as in Scholl's integrality condition.
\end{proof}

\begin{remark}
Formula \eqref{eq:invrelativelead} identifies the full leading
coefficient after the standard ideal class and unit normalizations
are inserted in $a_K/a_k$. The determinant statement is over
$\Q$ and therefore is unchanged by replacing a lattice by one of
finite index. The complex comparison has a kernel of dimension
$d_S$ before augmentation. This follows from \eqref{eq:sunitseq}
and the invertibility of $\alpha_{M_\varnothing}$.
\end{remark}

\subsection{Arbitrary critical lattices in products of norm tori}
\label{subsec:allcritical}

The choice of the ordinary unit lattice in Theorem~\ref{thm:relative}
can be removed. This gives a class of tori for which Scholl's
Conjecture~A holds for every critical constant lattice map. The split factor enters through its poles. Each copy of $\Gm$ supplies a
zeta pole which can cancel one of the local zeros.

Let $T_0=\prod_{i=1}^h T_{K_i/k_i}$, $T=T_0\times\Gm^t$, and
$M=[Y\to T]$ with $Y$ constant. Empty products are permitted.
Set $b=\sum_i b_0(K_i/k_i)$,
$a_0=\prod_i a_{K_i}/a_{k_i}$,
$D_q=\im\nu_q$, $r_q=\dim D_q$, and $r=\sum_q r_q$.
All these images are fixed by the local Galois group, since $Y$ is
constant. Write $[x]$ for the class of a nonzero real number modulo
$\Q^\times$.

\begin{theorem}[Critical value for an arbitrary lattice]\label{thm:allcritical}
Suppose that $\alpha_M$ is an isomorphism. Then $\dim Y_\Q=b+t$,
$r\ge t$, and
\begin{equation}\label{eq:arbitrary-leading}
 \ord_{s=0}L_{\rm inv}(M,s)=r-t,\qquad
 L_{\rm inv}^*(M,0)=(-1/2)^{b+t}a_0\prod_q(\log q)^{r_q}.
\end{equation}
If $r=t$, then
\[
 c^+(M)=\left[a_0\prod_q(\log q)^{r_q}\right].
\]
Thus the critical quotient is rational. It is nonzero exactly when
$r=t$. The rationality statement extends to tori isogenous to $T$
over $\Q$.
\end{theorem}

\begin{proof}
Proposition~\ref{prop:realdefect} and the ordinary unit computation
of Theorem~\ref{thm:relative} give $\dim C_T=b+t$. Criticality
identifies $Y_\R$ with $C_T$. Project to the $t$ real logarithms of the
split factor. For $D=\bigoplus_qD_q$ this surjection factors as
\begin{equation}\label{eq:splitprojection}
 Y_\R \xrightarrow{v=(\nu_q)} D_\R
 \xrightarrow{H_\infty} \R^t,\qquad
 H_\infty((x_q))=\sum_q(\log q)\pi_q(x_q).
\end{equation}
Here $\pi_q:D_q\to\Q^t$ is the projection on the split cocharacters.
The factorization follows from $|x|=\prod_q q^{\ord_q x}$ for
$x\in\Q^\times$. It gives $r\ge t$. Formula~\eqref{eq:globaleuler}
now reads
\[
 L_{\rm inv}(M,s)=\zeta(s)^{b+t}\zeta(s+1)^t
 \prod_i\frac{\zeta_{K_i}(s+1)}{\zeta_{k_i}(s+1)}
 \prod_q(1-q^{-s})^{r_q}.
\]
Taking the leading coefficients gives \eqref{eq:arbitrary-leading}.

If $r=t$, then $v$ is surjective over $\Q$, and $H_\infty$ is an
isomorphism. The rational map $(x_q)\mapsto\sum_q\pi_q(x_q)$ is an
isomorphism because $H_\infty$ is its composite with scalar
multiplication by $\log q$ on each summand. Let $Y_0=\ker v$. It has
dimension $b$. Its image in the split torus is torsion, and in each
$T_{K_i/k_i}$ it consists of ordinary relative units. Indeed the
vanishing of every local cocharacter valuation means that all finite
valuations of the corresponding field element vanish. The real
comparison restricted to $Y_0$ modulo torus periods is injective.
The rational image of $Y_0$ therefore equals the product of the
ordinary relative unit spaces, whose dimension is $b$.
In determinant bases adapted to $Y_0$ and to $T_0$, the first block
is the ordinary comparison for $T_0$ (including its compact Betti
periods), and the quotient block is $H_\infty$. Their determinant
classes are $[a_0]$ and $[\prod_q(\log q)^{r_q}]$, respectively.
For $r>t$ equation~\eqref{eq:arbitrary-leading} gives $L_{\rm inv}(M,0)=0$.
An isogeny transports the rational realizations, comparisons, and
the image of each valuation isomorphically.
\end{proof}

For $T=\Gm^t$ this theorem includes the split torus of every rank.
For $t=1$ and $1\mapsto q$, it gives Scholl's original example
\cite[\S II, example after Conjecture~A]{Scholl}. For a single relative
norm torus ($t=0$), a nonzero critical value forces the lattice to
span the ordinary relative units. The isogeny
$T_{K/\Q}\times\Gm\to\Res_{K/\Q}\Gm$ given by multiplication
places restrictions of scalars within the same result.

\section{The relative \texorpdfstring{$p$}{p}-adic residue formula}
\label{sec:padicrelative}

Fix an odd rational prime $p$ and an embedding $\Qbar\hookrightarrow
\Cp$. We use the Iwasawa logarithm, with $\log_p p=0$, on
the local multiplicative groups. For a finite extension $F/\Q_p$
write $H_f^1(F,\Q_p(1))$ for the Bloch--Kato finite subspace.
The local Kummer map and the inverse of the Bloch--Kato
exponential identify this subspace with the rationalized
completed unit group and then with $F$, respectively. In
these coordinates the inverse exponential on a unit class is
$\log_p$ \cite[\S3]{BlochKato}. The extension of $\log_p$ to
$F^\times$ uses the normalization on $p$. The valuation of
a general Kummer class measures its image in the quotient
by the finite subspace. For a restriction of scalars these
statements are applied at all the primes above $p$. The norm
becomes the trace in the logarithm coordinates since the
logarithm is additive on products of conjugates.

For a general Kummer class over $F$, the pair of maps
$(\log_p,v_p)$ identifies $H^1(F,\Q_p(1))$ with
$F\oplus\Q_p$, where $v_p(p)=1$. The Kummer class of
$p$ supplies a splitting of the second coordinate, and
the logarithm identifies the valuation-zero classes with
$F$. These extensions are semistable over $F$. The
valuation-zero part is crystalline, and the class of
$p$ is the extension in the Tate module of the Tate
curve of parameter $p$. Direct sums and subquotients
preserve semistability, so the same holds for the
extensions formed by adding and scaling these
classes. This is the split rank-one comparison behind
the norm logarithms in Section~\ref{sec:galois}.
Raynaud's geometric monodromy and its logarithmic Barsotti--Tate
realization treat the valuation part
\cite{Raynaud,BertapelleEtAl}. Over an unramified field, with a
constant lattice, split torus and good abelian quotient, the
semistable comparison in \cite[Thms.~3.4, 5.1]{MohajerSemistable}
identifies the valuation matrix with monodromy and the negative
logarithm of the integral part with the Hodge quotient of the
Frobenius lift. For a toric point $a$ over an unramified $K/\Q_p$
the two scalar characters used in Theorem~\ref{thm:cuplift} give
$[K:\Q_p]\ord_K(a)$ and $-\Tr_{K/\Q_p}\log_p(a)$.
Formula~\eqref{eq:local-sign} supplies the arithmetic pairing in
all the completions that occur below.

Suppose $K$ and $k$ are totally real. For
$M_\varnothing(K/k)$ define the local logarithm map
\begin{equation}\label{eq:relplog}
 \lambda_{K/k,p}:U_{K/k,\varnothing}\otimes_\Q\Q_p
 \longrightarrow
 \ker\left(\Tr_{K/k}:K\otimes_\Q\Q_p\to k\otimes_\Q\Q_p\right)
\end{equation}
by taking the logarithms in all completions of $K$ at $p$.
It is square by the unit rank calculation, and its determinant
in rational bases is well defined up to $\Q^\times$.
For a totally real field $L$, let $\zeta_{L,p}$ be the
Deligne--Ribet zeta function on the branch interpolating
the values of $\zeta_L^{(p)}(1-n)$ for positive
$n\equiv0\bmod(p-1)$. The superscript $(p)$ means that all Euler
factors above $p$ are removed. We use the interpolation of Deligne
and Ribet and the independent construction of Cassou-Nogu\`es
\cite{DR,Cassou}. Put
\[
 E_p(L)=\prod_{\mathfrak v\mid p}(1-(N\mathfrak v)^{-1}),
 \qquad
 a_{L,p}=\left.((s-1)\zeta_{L,p}(s))\right|_{s=1}.
\]
The second expression is the residue, including the value zero
when the pole cancels.

\begin{theorem}\label{thm:prelative}
Assume $a_{k,p}\ne0$. The quotient
$\zeta_{K,p}(s)/\zeta_{k,p}(s)$ is analytic near $1$, and
\begin{equation}\label{eq:relativepadicvalue}
 \left.\frac{\zeta_{K,p}(s)}{\zeta_{k,p}(s)}\right|_{s=1}
 \frac{E_p(k)}{E_p(K)}
 \ \in\ \Q^\times\det(\lambda_{K/k,p}).
\end{equation}
In particular the normalized value is nonzero if and only if
$K$ satisfies Leopoldt's conjecture at $p$. When both fields
are abelian over $\Q$, the hypothesis and the nonvanishing
hold unconditionally.
\end{theorem}

\begin{proof}
For the absolute unit group of a totally real $L$ append
the trace-one column $1/[L:\Q]$ to its logarithm map into
$L\otimes\Q_p$. The embedding matrix of an integral basis
has determinant a choice of $\sqrt{\operatorname{disc}L}$.
Replacing one row by the sum of the rows shows that the
absolute determinant represents
$R_{L,p}/\sqrt{\operatorname{disc}L}$ in rational bases, where
$R_{L,p}$ is the $p$-adic regulator. The trace of all unit
logarithms is zero because the norm of an ordinary unit to
$\Q$ is $\pm1$. As in \eqref{eq:absaug}, the norm and trace
exact sequences identify the relative determinant with
the ratio of the absolute determinants when the denominator
is nonzero. This determinant computation is valid if the
numerator vanishes.

Colmez's $p$-adic class number formula
\cite[pp.~371--372]{Colmez} gives, with coherent orientations,
\[
 a_{L,p}=E_p(L)
 \frac{2^{[L:\Q]}h_L R_{L,p}}
      {w_L\sqrt{\operatorname{disc}L}},
\]
where $h_L$ is the class number and $w_L$ the number of roots
of unity. Since $a_{k,p}\ne0$, the analytic germs
$(s-1)\zeta_{K,p}(s)$ and $(s-1)\zeta_{k,p}(s)$ have a
quotient with value $a_{K,p}/a_{k,p}$. The determinant
identity proves \eqref{eq:relativepadicvalue}. The
nonvanishing of $R_{L,p}$ is the regulator formulation
of Leopoldt's conjecture. Brumer proves this conjecture
for abelian extensions of $\Q$ \cite{Brumer}.
\end{proof}

\subsection{Adding primes away from \texorpdfstring{$p$}{p}}

Let $K/k$ be totally real, assume $a_{k,p}\ne0$, and
let $S$ be a finite set of rational primes disjoint from
$\{p\}$. For a $p$-adic unit $q$, $\langle q\rangle$
denotes its projection to $1+p\Z_p$. Consider the
analytic function near zero given by
\begin{equation}\label{eq:padicS}
 \mathcal Z_{K/k,S,p}(s)=
 \frac{\zeta_{K,p}(1+s)}{\zeta_{k,p}(1+s)}
 \prod_{q\in S}(1-\langle q\rangle^{-s})^{d_q}.
\end{equation}
This is the $p$-adic interpolation of the torus
quotient and finite correction in \eqref{eq:invrelative},
with the lattice zeta factor divided out. More precisely,
at each interpolation point $s=-n$ with $n>0$ divisible by $p-1$,
the quotient in \eqref{eq:padicS}, wherever defined, becomes the ratio of
the $p$-depleted complex zeta values at $1-n$ and the
last factors become $(1-q^n)^{d_q}$.

Define $\lambda_{K/k,S,p}$ on $U_{K/k,S}\otimes\Q_p$
using the local logarithm into $\Lie(T_{K/k})\otimes\Q_p$
and the additional rows $(\log_p q)\ord_q$ for $q\in S$.
Its target is
\[
  \bigl(\Lie(T_{K/k})\otimes\Q_p\bigr)
   \oplus\bigoplus_{q\in S}\bigl(D_q(K/k)\otimes\Q_p\bigr).
\]
All points in the lattice are units at the places above
$p$, so the logarithms belong to the finite Kummer comparison.

\begin{corollary}\label{cor:padicS}
The order at zero of \eqref{eq:padicS} is at least $d_S$, and
\[
 \frac{E_p(k)}{E_p(K)}[s^{d_S}]\mathcal Z_{K/k,S,p}(s)
       \ \in\ \Q^\times\det(\lambda_{K/k,S,p}).
\]
The order is exactly $d_S$ if and only if $K$
satisfies Leopoldt's conjecture at $p$. In particular
this equality holds for $K/\Q$ abelian when $k\subset K$.
\end{corollary}

\begin{proof}
Use \eqref{eq:sunitseq} to take bases adapted to the ordinary
units. The diagonal blocks of the comparison are the
ordinary relative logarithm and
$\bigoplus_{q\in S}(\log_p q)\mathrm{id}_{D_q(K/k)}$.
The logarithm of a rational prime $q\ne p$ is nonzero
because a rational $p$-adic unit in the kernel is torsion.
The expansion $1-\langle q\rangle^{-s}=s\log_p q+O(s^2)$
now gives the formula and the order statement from
Theorem~\ref{thm:prelative}.
\end{proof}

\subsection{Cancellation of the split poles at \texorpdfstring{$p$}{p}}

Return to the notation of Theorem~\ref{thm:allcritical}. Assume
that all $k_i,K_i$ are totally real, $a_{k_i,p}\ne0$, $M$ is
critical at the real place, and $r_p=0$. The last condition means
that the lattice points are units at every place over $p$.
In the cyclotomic variable centered at zero put
\begin{equation}\label{eq:splitpadicgerm}
 Z_{M,p}(s)=\zeta_{\Q,p}(1+s)^t
 \prod_i\frac{\zeta_{K_i,p}(1+s)}{\zeta_{k_i,p}(1+s)}
 \prod_{\substack{q\ne p\\r_q>0}}(1-\langle q\rangle^{-s})^{r_q}.
\end{equation}
The shifted zeta factors interpolate the toric part. As in
\eqref{eq:padicS}, the weight-zero lattice factor has been divided
out. Write $\lambda_{M,p}:Y_{\Q_p}\to\Lie(T)_{\Q_p}$ for the
componentwise local logarithm in rational bases. This map is square
because all the fields are totally real.

\begin{corollary}\label{cor:splitpadic}
The function \eqref{eq:splitpadicgerm} is analytic at zero, with
order at least $r-t$. Equality holds precisely when every $K_i$
satisfies Leopoldt's conjecture at $p$. If $r=t$, then
\[
 E_p(\Q)^{-t}\prod_i\frac{E_p(k_i)}{E_p(K_i)}\, Z_{M,p}(0)
       \ \in\ \Q^\times\det\lambda_{M,p}.
\]
Here $\Q^\times\cdot0=\{0\}$ if the logarithm map is singular.
\end{corollary}

\begin{proof}
Expand $(s\zeta_{\Q,p}(1+s))^t$ at zero. Its value is
$E_p(\Q)^t$ by Colmez's formula. Each finite factor has first
coefficient $\log_p q\ne0$, and $r\ge t$ by
Theorem~\ref{thm:allcritical}. Theorem~\ref{thm:prelative}
gives the order assertion and identifies the normalized coefficient
of $s^{r-t}$ up to $\Q^\times$ with
$\prod_i\det\lambda_{K_i/k_i,p}\prod_q(\log_p q)^{r_q}$.
For $r=t$, use the same rational spaces $Y_0$ and $D$ as in
\eqref{eq:splitprojection}. On a rational split coordinate $x$,
$\log_p x=\sum_{q\ne p}\ord_q(x)\log_p q$. The quotient block of
the local comparison is therefore $(x_q)\mapsto\sum_q(\log_p q)
\pi_q(x_q)$, whose determinant differs from $\prod_q(\log_p q)^{r_q}$
by a nonzero rational factor. On $Y_0$ the blocks are the ordinary
relative logarithm maps. Multiplication of these determinants proves
the formula, even when one of the ordinary blocks is singular.
\end{proof}

\subsection{An integrality criterion in rank one}

Let $K$ be a real quadratic field. Write $T=T_{K/\Q}$,
$M_a=[\Z\xrightarrow{a}T]$ for a norm-one point
$a\in K^\times$ of infinite order, and choose a generator
$\varepsilon$ of the norm-one unit group modulo torsion.
Fix a real embedding $\sigma:K\hookrightarrow\R$, and
write $\iota$ for the chosen $p$-adic embedding on $K$.
The archimedean comparison is an isomorphism and
\[
 c^+(M_a)=\left[\log|\sigma(a)|/\sqrt{\operatorname{disc}K}\right].
\]
Indeed the torus is split over $\R$, so the plus
comparison is one logarithm. Passing from an
embedding coordinate to a rational trace-zero basis
divides by $\sqrt{\operatorname{disc}K}$ up to $\Q^\times$.
The $p$-adic logarithm of $M_a$ in the same rational
trace-zero basis represents
$c_p(M_a)= [\log_p(\iota(a))/\iota(\sqrt{\operatorname{disc}K})]$.
This is the finite Kummer period when $a$ is a local unit
above $p$ and denotes the specified branch logarithm in general.

\begin{theorem}\label{thm:rigidity}
For any such $a$, the following conditions are equivalent.
\begin{enumerate}
\item $a$ is a global unit.
\item $L_{\gr}(M_a,0)/c^+(M_a)$ is algebraic.
\item For the relative $p$-adic value in
\eqref{eq:relativepadicvalue} for $K/\Q$, the quotient
of that value by $c_p(M_a)$ is algebraic over $\Q$.
\end{enumerate}
When these conditions hold, both quotients are nonzero
rational numbers for any choices of rational bases.
\end{theorem}

\begin{proof}
First $\log_p(\iota(a))\ne0$. An element $x$ of a finite
extension of $\Q_p$ with $\log_p x=0$ has a power of the
form $p^m$ for some integer $m$. To see this, raise $x$
to a power to make its valuation integral, divide by
the corresponding power of $p$, and use that the kernel
of the logarithm on local units is torsion. For $a$ the
resulting identity is algebraic and can be normed to
$\Q$. The norm-one condition forces $m=0$, which
contradicts the assumption of infinite order.

The graded complex $L$-function is the same for $M_a$ and
$M_\varepsilon$. Theorems~\ref{thm:relative} and
\ref{thm:prelative} reduce the algebraicity assertions to
the assertions that
\[
 \frac{\log|\sigma(\varepsilon)|}{\log|\sigma(a)|},
 \qquad
 \frac{\log_p(\iota(\varepsilon))}{\log_p(\iota(a))}
\]
are algebraic, respectively. If $a$ is a unit, it is
equal to a power of $\varepsilon$ up to torsion, and
both ratios are rational. If $a$ is not a unit, some
finite valuation of $a$ is nonzero. A rational relation
between the real logarithms would give $a^m\varepsilon^n=\pm1$
by injectivity of $\sigma$. The finite valuation forces
$m=0$, and then $n=0$. A rational relation between
the $p$-adic logarithms likewise gives a norm-one element
of logarithm zero, hence a torsion element by the first
argument. The same valuation proves rational independence.
Baker's theorem for real logarithms and its $p$-adic
analogue due to Brumer give algebraic linear independence
\cite{Baker,Brumer}. In each realization the ratio is
therefore transcendental if $a$ is not a unit.
\end{proof}

\begin{example}\label{ex:quadratic}
Take $K=\Q(\sqrt5)$ and $a=(4+\sqrt5)/(4-\sqrt5)$.
The numerator and denominator have norm $11$ and generate
the two distinct primes at $11$. Thus $a$ has norm one
and both algebraic-quotient conditions fail. At $11$ the
valuation map of $M_a$ has rank one. Formula
\eqref{eq:globaleuler} gives
$\ord_{s=0}L_{\mathrm{inv}}(M_a,s)=1$ although
$L_{\gr}(M_a,0)\ne0$. The criticality of $M_a$ follows
from its nonzero real logarithm. The full relative
$\{11\}$-unit lattice contains an ordinary unit direction
as well as the valuation direction. Theorem~\ref{thm:relative}
computes its augmented determinant.
\end{example}

\section{The Gross regulator on a Galois CM field}
\label{sec:galois}

Let $N/\Q$ be a finite Galois CM field, $G=\Gal(N/\Q)$,
and $j\in G$ its central complex conjugation. Write
$S_p(N)$ for the places above $p$ and use the rational
$G$-modules
\[
 U=\bigl(\OO_N[1/p]^\times\otimes_\Z\Q\bigr)^-,
 \qquad P=\Q[S_p(N)]^- ,
\]
where minus means the $(-1)$-eigenspace for $j$.
Any finite family of CM fields admits such a common overfield.
Take the compositum of all their conjugates. Complex conjugation
restricts to the canonical CM involution on each conjugate
field, so it is central in the Galois group of the compositum.
Its fixed field is totally real and has index two.
If $u$ is an ordinary unit, $u/j(u)$ has absolute value
one at every complex embedding and is an algebraic integer,
so it is torsion by Kronecker's theorem. It follows that
the odd part of the rational ordinary units vanishes.
Finiteness of the class group now implies that the
valuation map from $U$ to $P$ is an isomorphism.

The relevant 1-motive is $M_{\{p\}}(N/N^+)$, where $N^+=N^{\langle
j\rangle}$. Its rational lattice is $U$ and its torus is the
kernel of $N_{N/N^+}$. An element of $G$ acts both on the
constant lattice of rational points and as an automorphism
of that torus over $\Q$. In particular the $G$-action
gives genuine endomorphisms of the 1-motive. The action on
the lattice here is by endomorphisms of the motive.
The lattice has trivial absolute Galois action as a group
scheme over $\Q$. The odd valuation and norm logarithm maps
are obtained by evaluating the Kummer classes of this lattice
under the local characters at the places above $p$.

\subsection{Valuations of local norms}

Let $f_N$ be the common residue degree of the primes of
$N$ above $p$. Define two $G$-equivariant maps
\begin{align}
 O_N:U&\longrightarrow P,&
       O_N(u)_{\mathfrak w}
       &= \ord_p N_{N_\mathfrak w/\Q_p}(u)
        = f_N\ord_\mathfrak w(u),\label{eq:ON}\\
 \Lambda_N:U&\longrightarrow P\otimes_\Q\Q_p,&
       \Lambda_N(u)_{\mathfrak w}
       &= -\log_p N_{N_\mathfrak w/\Q_p}(u).\label{eq:LN}
\end{align}
The minus sign in the second map is the convention of
Dasgupta--Kakde--Ventullo. By the valuation isomorphism
above, the endomorphism
\begin{equation}\label{eq:RN}
 R_N=\Lambda_NO_N^{-1}\in\End_{\Q_p[G]}(P\otimes_\Q\Q_p)
\end{equation}
is defined. If the logarithm on $\Q_p^\times$ is changed
to $\log_p+c\,\ord_p$, the new endomorphism is $R_N-cI$.
This follows by subtracting $cO_N$ in \eqref{eq:LN}.
We use $\log_p p=0$ below.

\subsection{Evaluation on the endomorphism algebra}

We identify $\Qbar$ with its image in $\Cp$ and set
$P_{\mathrm a}=P\otimes_\Q\Qbar$. Choose a place
$\mathfrak P$ above $p$ through the fixed embedding of $N$
into $\Cp$, and put $D=G_{\mathfrak P}$. Use the idempotents
\[
 e_D=|D|^{-1}\sum_{d\in D}d,\qquad e^-=(1-j)/2,\qquad
 e=e_D e^- \quad\text{in }\Qbar[G].
\]
The identification of $S_p(N)$ with $G/D$ gives
$P_{\mathrm a}=\Qbar[G]e$. If $e=0$, all the
endomorphism and coefficient spaces below are zero,
with the usual empty determinant equal to $1$.

\begin{lemma}\label{lem:log}
The map $U\otimes\Qbar\to\Cp$ given by $u\mapsto
\log_p(\iota(u))$ is injective. Its image has zero
intersection with $\Qbar$.
\end{lemma}

\begin{proof}
On $N^\times\otimes_\Z\Q$ the kernel of the single-embedding
logarithm is the line spanned by $p$. The argument with local
units in the proof of Theorem~\ref{thm:rigidity} gives
this assertion after clearing denominators in any element
of the rationalized multiplicative group. The line is even
under $j$, so the rational logarithm is injective on $U$.
Baker--Brumer therefore gives the asserted injectivity
over $\Qbar$ \cite{Brumer,Maksoud}.

For the last assertion, choose representatives
$u_1,\ldots,u_d$ of a rational basis of $U$ in the
norm-one $p$-units. By taking algebraic powers with a
positive integer exponent and dividing by powers of $p$,
one can replace their images by algebraic local units
$z_i$ whose logarithms are nonzero rational multiples
of the original ones. Further positive powers put all
$z_i$ in the domain of the local exponential. The
logarithms remain rationally independent. Suppose that
$1,\log_p z_1,\ldots,\log_p z_d$ satisfy an algebraic
linear equation. The $p$-adic analytic subgroup theorem
applied to the algebraic point $(1,z_1,\ldots,z_d)$ in
$\Ga\times\Gm^d$ gives an algebraic subgroup whose Lie
algebra is contained in the corresponding proper
hyperplane and contains the logarithm vector
$(1,\log_p z_1,\ldots,\log_p z_d)$
\cite[Th.~1.1 and \S3.1]{FuchsPham}. A connected subgroup of
$\Ga\times\Gm^d$ is the product of its unipotent and
torus parts. This follows from Jordan decomposition in a
faithful linear representation of the commutative group.
The first logarithm coordinate forces the unipotent
part to be $\Ga$. A proper torus part would impose a
nonzero integral character relation on the torus
logarithms. Rational independence rules this out. Thus
the Lie algebra cannot lie in a proper hyperplane.
\end{proof}

\begin{proposition}\label{prop:independent-end}
Let $\mathcal E=\End_{\Qbar[G]}(P_{\mathrm a})$.
View $R_N$ in $\mathcal E\otimes_{\Qbar}\Cp$ and extend algebraic
linear functionals to this scalar extension. The map
\begin{equation}\label{eq:evaluation}
 \Qbar\oplus\mathcal E^\vee\longrightarrow\Cp,\qquad
 (b,\varphi)\longmapsto b+\varphi(R_N)
\end{equation}
is injective. Equivalently, the coordinates of $R_N$ in
any $\Qbar$-basis of $\mathcal E$ are linearly independent
over $\Qbar$ together with $1$.
\end{proposition}

\begin{proof}
Let $\mathcal A=e\Qbar[G]e$. An equivariant endomorphism
of $\Qbar[G]e$ is right multiplication by a unique
$a\in\mathcal A$, determined by the image of $e$.
For $x\in\mathcal A$, the functional on these
endomorphisms given by the $\mathfrak P$ coordinate of
$xa$ is $a\mapsto |D|[1](xa)$, where $[1]$ means
the identity coefficient in the group algebra.
The pairing $(x,a)\mapsto[1](xa)$ is nondegenerate.
In the full group algebra this follows by pairing a
group element with its inverse. If $a=eae$ then
$[1](za)=[1](eze\,a)$ by cyclicity of the identity
coefficient. Thus compressing both arguments to
$\mathcal A$ preserves nondegeneracy. It follows that
every $\varphi\in\mathcal E^\vee$ is evaluation of the
$\mathfrak P$ coordinate on a unique $x\in\mathcal A$.

For this $x$ take $u=O_N^{-1}x$. Since $e_D x=x$,
the vector $u$ is $D$-invariant. The norm on the
completion at $\mathfrak P$ is the product over $D$.
Extending linearly to algebraic coefficients thus gives
\[
 \varphi(R_N)=(R_Nx)_{\mathfrak P}
  =-|D|\log_p(\iota(u)).
\]
The map $x\mapsto u$ is injective and
Lemma~\ref{lem:log} gives \eqref{eq:evaluation}.
\end{proof}

\subsection{Restriction to character spaces}

Let $F\subset N^+$ be a totally real subfield,
$B=\Gal(N/F)$, and let $\chi:B\to\Qbar^\times$ be a
one-dimensional character with $\chi(j)=-1$. Put
$H=N^{\ker\chi}$ and let $e_{\chi^{-1}}\in\Qbar[B]$
be the projector onto the $\chi^{-1}$-eigenspace.
The extension $H/F$ is cyclic and CM. Let
\[
 \mathcal R(F,\chi)=\{\mathfrak v\mid p\text{ in }F:
                      \chi|_{G_{\mathfrak v}}=1\},\qquad
 r(F,\chi)=|\mathcal R(F,\chi)| .
\]
Here $G_\mathfrak v$ is a decomposition group in $B$.
Because $H$ is cut out by $\chi$, at each
$\mathfrak v\in\mathcal R(F,\chi)$ the chosen completion
of $H$ is canonically $F_\mathfrak v$. Choose one prime
of $H$ above each $\mathfrak v\in\mathcal R(F,\chi)$.
Write $f_\mathfrak v=f(F_\mathfrak v/\Q_p)$.
On the $\chi^{-1}$ part of the rational $p$-units of
$H$ with algebraic coefficients, define the matrices
$O_\chi$ and $L_\chi$ using the rows
\[
 \ord_\mathfrak v,\qquad
 -\log_p\circ N_{F_\mathfrak v/\Q_p},
 \qquad \mathfrak v\in\mathcal R(F,\chi).
\]
The columns refer to a common basis of the character space
over the embedded algebraic character field. Any later $p$-adic
coefficient field is taken large enough to contain that basis. The valuation matrix is invertible.
Indeed, the divisor permutation module of $H$ has
one $\chi^{-1}$ direction over $\mathfrak v$ if its
decomposition group acts trivially under $\chi$, and zero
otherwise. The ordinary units have no odd part. Define
\begin{equation}\label{eq:normalizedmatrix}
 B(F,\chi)=L_\chi O_\chi^{-1}\operatorname{diag}(f_\mathfrak v^{-1}) .
\end{equation}
Changing the chosen primes of $H$ conjugates this matrix by
a diagonal algebraic matrix. The matrix $L_\chi O_\chi^{-1}$
before the last diagonal factor has determinant equal to the
Gross regulator in the convention of \cite[\S1.1]{DKV}.

\begin{proposition}\label{prop:descent}
The operator $B(F,\chi)$ is algebraically conjugate to
the restriction of $R_N$ to
$(e_{\chi^{-1}}P_{\mathrm a})\otimes_{\Qbar}\Cp$. The conjugating
matrix has algebraic entries.
Consequently its restricted operator and its characteristic
polynomial can be transported to any common Galois CM
overfield using the same normalization.
\end{proposition}

\begin{proof}
Taking $\ker\chi$-invariants in the rational $p$-units
of $N$ gives the rational $p$-units of $H$. This follows
by averaging with the norm to $H$ and division by
$[N:H]$. The same assertion on the minus character
space then holds over $\Qbar$. On
$e_{\chi^{-1}}P_{\mathrm a}$, a coordinate at one
prime of $N$ over each $\mathfrak v\in\mathcal R(F,\chi)$
gives a basis of coordinate functionals. Equivariance
determines the other coordinates from the selected
ones. A prime orbit whose stabilizer has a nontrivial
$\chi$ value has zero coordinates.

At the selected prime, write $b_\mathfrak v=[N_\mathfrak w:
F_\mathfrak v]$. For a vector of $p$-units from $H$, its
valuation-of-norm and logarithm-of-norm in $N_\mathfrak w$
both multiply by $b_\mathfrak v$, by transitivity of the
norm. The valuation of the norm from $F_\mathfrak v$
itself is $f_\mathfrak v\ord_\mathfrak v$. Thus the two
matrices for $O_N$ and $\Lambda_N$ in the selected
coordinate rows are
\[
 \operatorname{diag}(b_\mathfrak v)
       \operatorname{diag}(f_\mathfrak v)O_\chi,
 \qquad
 \operatorname{diag}(b_\mathfrak v)L_\chi .
\]
Their quotient is conjugation of \eqref{eq:normalizedmatrix}
by $\operatorname{diag}(b_\mathfrak v)$. The same calculation
proves the assertion upon further enlargement of $N$.
\end{proof}

\subsection{A local lifting criterion for the Gross matrix}
\label{subsec:reciprocity}

We now give the matrix in \eqref{eq:normalizedmatrix} a local
cohomological interpretation. This identifies the local input to the
global special-value formula and gives a meaning to every principal
minor. Write $\mathcal R=\{\mathfrak v_1,\ldots,\mathfrak v_r\}$,
$K_i=F_{\mathfrak v_i}$, and let $E$ be a finite extension of $\Q_p$
containing the values of $\chi$. For a $p$-adic local field $K$ let
$\kappa_K$ denote its Kummer map and let $\operatorname{inv}_K$ be
the invariant in degree two. Identify additive continuous characters
$\eta:K^\times\to E$ with $H^1(K,E)$ by the reciprocity convention
\begin{equation}\label{eq:local-sign}
 \operatorname{inv}_K(\kappa_K(x)\cup\eta)=\eta(x)
 \qquad(x\in K^\times).
\end{equation}
The character is in the second cup factor. This is local Tate
duality \cite[Ch.~I, \S2]{Milne}, in the sign convention of
\cite[eq.~(2.3)]{MohajerSemistable}. At $\mathfrak v_i$ set
\[
 o_i=\ord_p\circ N_{K_i/\Q_p},\qquad
 \ell_i=-\log_p\circ N_{K_i/\Q_p}.
\]
The two characters are independent because restriction to
$\Q_p^\times\subset K_i^\times$ gives nonzero common degree
multiples of the valuation and negative logarithm.

To specify the extension whose dual we use, take
$T_{\mathcal R}=\prod_i\Res_{K_i/\Q_p}\Gm$ and map the rational
odd $p$-units of $H$ to $T_{\mathcal R}(\Q_p)\otimes\Q$ by the
chosen local embeddings. A lattice up to isogeny defines the
corresponding local toric 1-motive. Extend its Tate realization to
$E$ and pull back the lattice quotient to the $\chi^{-1}$-subspace
$Y_\chi$. Here $G_{\Q_p}$ denotes the absolute Galois group of $\Q_p$.
The pullback gives an exact sequence of $E[G_{\Q_p}]$-modules
\[
 0\longrightarrow V_p(T_{\mathcal R})_E
 \longrightarrow V_\chi\longrightarrow Y_\chi\longrightarrow0
\]
with trivial local Galois action on $Y_\chi$. Its extension classes
are the local Kummer classes of the chosen embeddings of the
lattice points. For an $E$-linear representation write $V^\dagger=V^\vee(1)$,
with the $E$-linear dual.
By Shapiro's lemma $H^1(\Q_p,V_p(T_{\mathcal R})_E^\dagger)$ is
$\bigoplus_i H^1(K_i,E)$. A tuple in this sum \emph{lifts} if it is
in the image of $H^1(\Q_p,V_\chi^\dagger)$.

\begin{theorem}[Gross matrices and dual lifts]\label{thm:cuplift}
Put $B=B(F,\chi)$. For columns $x,y\in E^r$, the tuple
$(x_i o_i+y_i\ell_i)_i$ lifts if and only if
\begin{equation}\label{eq:liftgraph}
                         x=-B^{\mathsf t}y .
\end{equation}
In particular, for $\mathbf t=(t_1,\dots,t_r)$ the space of lifting
tuples of the form $(y_i(\ell_i-t_i o_i))_i$ has dimension
$\dim\ker(B^{\mathsf t}-\operatorname{diag}(\mathbf t))$.
The polynomial defining the determinant condition for a nonzero lift is
\begin{equation}\label{eq:principal-cup}
 \det(\operatorname{diag}(\mathbf t)-B)=
 \sum_{J\subseteq\{1,\ldots,r\}}(-1)^{|J|}
 \det B_{J,J}\prod_{i\notin J}t_i .
\end{equation}
The statements hold after any finite extension of $E$.
\end{theorem}

\begin{proof}
Local invariants identify $H^2(\Q_p,Y_\chi^\vee(1))$ with
$Y_\chi^\vee$. The connecting map of the dual sequence, up to the
overall sign of the dual extension, takes a tuple $(\eta_i)$ to
\[
 u\longmapsto \sum_i\operatorname{inv}_{K_i}
       (\kappa_{K_i}(u_{\mathfrak v_i})\cup\eta_i).
\]
Indeed, dualizing the extension transposes its Kummer cocycle with a
minus sign. Cup product with the dual class evaluates this cocycle on
$u$. Shapiro's lemma takes the resulting pairing over $\Q_p$ to the
sum of the corestricted local cup products. Corestriction preserves
the value of the local invariant.
This is the adjoint Kummer lifting criterion of
\cite[Prop.~7.5, eq.~(7.8)]{MohajerSemistable}. The proof from
\eqref{eq:local-sign} works for ramified $K_i$ as well.

In the column basis for $Y_\chi$ used in \eqref{eq:normalizedmatrix},
let $\widetilde O=\operatorname{diag}(f_{\mathfrak v_i})O_\chi$.
By \eqref{eq:local-sign}, the row of the displayed functional equals
$x^{\mathsf t}\widetilde O+y^{\mathsf t}L_\chi$.
Exactness says that the tuple lifts exactly when this row is zero.
Right multiplication by $\widetilde O^{-1}$ gives
\eqref{eq:liftgraph}. Substituting $x=-\operatorname{diag}(\mathbf t)y$
proves the dimension formula. Expansion of the determinant by the
chosen diagonal variables gives \eqref{eq:principal-cup}, with the
empty minor equal to one.
\end{proof}

With all $t_i=t$, \eqref{eq:principal-cup} is the characteristic
polynomial of $B$. With all $t_i=0$, its constant term is
$(-1)^r$ times the normalized Gross--Stark coefficient of
Theorem~\ref{thm:analytic}. In particular exact order is equivalent
to the absence of a nonzero lifting tuple made entirely of the
$\ell_i$. For $r=1$ the unique local slope $t$ at which a lift
appears is that normalized first derivative. For any $J$ the minor
in \eqref{eq:principal-cup} is the corresponding minor of the
unscaled Gross matrix divided by $\prod_{i\in J} f_{\mathfrak v_i}$.

Dasgupta and Spiess formulate a different description of the principal
minors using a global Eisenstein cocycle \cite{DS}. Honnor compares
two formulas for the diagonal terms over a totally real cubic
field \cite{Honnor}. Dasgupta and Kakde prove an integral
reciprocity refinement for Brumer--Stark units \cite{DKIntegral}.
Dasgupta, Honnor and Spiess identify the multiplicative-integral and
two cohomological formulas for those units \cite{DHS}. The local formula above uses the actual Kummer
classes and records exactly which dual characters lift. In
\cite[Thms.~13.3, 13.4]{MohajerSemistable} one finds explicit Herr
cochains for the toric basis family and relative lifts under the
coefficient hypotheses given there. The division-tower comparison in
\cite[Thms.~15.6, 15.7 and \S16.2]{MohajerSemistable} identifies a
period pairing with an arithmetic Iwasawa pairing. The analytic
function used here comes from Deligne--Ribet interpolation, and its
connection with the determinant is the global theorem of
Dasgupta, Kakde and Ventullo.

\section{Joint relations and analytic Taylor coefficients}
\label{sec:joint}

For an irreducible $\Qbar$-representation $V_\rho$ of $G$ with
$j$ acting by $-1$, put
\[
  W_\rho=\Hom_G(V_\rho,P_{\mathrm a}),\qquad
  d_\rho=\dim W_\rho=\dim V_\rho^{D}.
\]
The dimension identity follows by Frobenius reciprocity for
the permutation module on $G/D$. In the evaluation
decomposition
\begin{equation}\label{eq:decomp}
 P_{\mathrm a}=\bigoplus_{\rho\text{ odd}} V_\rho\otimes W_\rho,
 \qquad
 R_N=\bigoplus_{\rho\text{ odd}} 1_{V_\rho}\otimes C_\rho ,
\end{equation}
the formula for $R_N$ is taken after extending scalars to $\Cp$.
Choose algebraic bases of $W_\rho$. By
Proposition~\ref{prop:independent-end}, all the entries of all
$C_\rho$ are linearly independent over $\Qbar$ together
with $1$. In fact the matrix coordinate functionals form
a basis for the dual of
$\mathcal E=\prod_{\rho\ {\rm odd}}\End_{\Qbar}(W_\rho)$.
Terms with $d_\rho=0$ can be omitted.

For a pair $(F,\chi)$ as above put
\[
 m_{F,\chi}(\rho)=\dim_{\Qbar}e_{\chi^{-1}}V_\rho .
\]
This equals the multiplicity of $\rho$ in
$\operatorname{Ind}_{\Gal(N/F)}^G\chi^{-1}$. Restricting
\eqref{eq:decomp} shows that
$r(F,\chi)=\sum_\rho m_{F,\chi}(\rho)d_\rho$.

\begin{theorem}\label{thm:joint}
Fix a finite family of pairs $(F_i,\chi_i)$ in the same
Galois CM field $N$. For $B_i=B(F_i,\chi_i)$ there are
simultaneous algebraic descriptions
\begin{equation}\label{eq:blockpoly}
 B_i\ \sim\
 \bigoplus_{\rho\text{ odd}} C_\rho^{\oplus m_{F_i,\chi_i}(\rho)},
 \qquad
 \det(TI-B_i)=
 \prod_{\rho\text{ odd}}\det(TI-C_\rho)^{m_{F_i,\chi_i}(\rho)} .
\end{equation}
In particular, the $\Qbar$-dimension of the combined span
of the entries of all $B_i$ equals
\begin{equation}\label{eq:lineardimension}
 \sum_{\substack{\rho\ {\rm odd}\\
                 m_{F_i,\chi_i}(\rho)>0\text{ for some }i}}
                 d_\rho^2 .
\end{equation}
Every algebraic linear relation among the entries is obtained
by substituting the repeated blocks of \eqref{eq:blockpoly}
after the algebraic basis changes in that equation. Adjoining
$1$ increases the span dimension by one.
\end{theorem}

\begin{proof}
Apply Proposition~\ref{prop:descent} to each pair and extend
\eqref{eq:decomp} along the projector $e_{\chi_i^{-1}}$.
The resulting block on
$e_{\chi_i^{-1}}V_\rho\otimes W_\rho$ is
$1\otimes C_\rho$, and choosing bases gives the specified
algebraic conjugation. Conjugation by an algebraic matrix
preserves the algebraic span of all entries. In block
coordinates this span has precisely the independent
coordinate functions indexed in \eqref{eq:lineardimension}.
Proposition~\ref{prop:independent-end} implies that a linear
combination of these functions vanishes at $R_N$ only
when it vanishes identically. The same argument includes
an independent constant term.
\end{proof}

\begin{corollary}\label{cor:traces}
For pairs $(F_i,\chi_i)$ as above and
$a_i,b\in\Qbar$, one has
\[
 b+\sum_i a_i\Tr B_i=0
 \quad\Longleftrightarrow\quad
 b=0\ \ \text{and}\ \
 \sum_i a_i m_{F_i,\chi_i}(\rho)=0
 \quad\text{for every odd }\rho\text{ with }d_\rho>0 .
\]
\end{corollary}

\begin{proof}
By \eqref{eq:blockpoly}, $\Tr B_i$ is the corresponding
linear combination of $\Tr C_\rho$. Each $\Tr C_\rho$
uses disjoint nonempty sets of the independent matrix
coordinates, so the traces and $1$ are independent.
\end{proof}

\begin{corollary}\label{cor:liftbound}
For $(F,\chi)$ let $\mathcal L_{\rm loc}$ be the space of tuples
$(y_i\ell_i)_i$ which lift in Theorem~\ref{thm:cuplift}. Then
\[
 \dim \mathcal L_{\rm loc}
 =\sum_\rho m_{F,\chi}(\rho)\dim\ker C_\rho
 \ \le\ \sum_\rho m_{F,\chi}(\rho)\lfloor d_\rho/2\rfloor .
\]
\end{corollary}

\begin{proof}
Take $\mathbf t=0$ in the local lifting criterion and use
\eqref{eq:blockpoly}. For an odd irreducible $\rho$ the rank
bound of Maksoud \cite[Th.~4.3 and the proof of Th.~4.6]{Maksoud}
gives $\dim\ker C_\rho\le\lfloor d_\rho/2\rfloor$.
Indeed his odd regulator map is the norm-logarithm map on the
multiplicity space for the odd $p$-unit module. Composing it with
the inverse valuation isomorphism preserves its rank. Replacing
$\rho$ by its contragredient in the character convention preserves
the local fixed-space dimension.
\end{proof}

\subsection{Gross--Stark coefficients}

Write $L_F(\chi,s)$ for the complex Artin $L$-function in
the Hecke character convention of \cite{DKV}. At a prime
$\mathfrak v\mid p$ set $\chi(\mathfrak v)=0$ if $\chi$
is ramified there, with the usual unramified convention
otherwise. Thus $\mathcal R=\mathcal R(F,\chi)$ is
exactly the subset where $\chi(\mathfrak v)=1$.
The $p$-depleted complex function
$L_F^{(p)}(\chi,s)$ has order $r=r(F,\chi)$ at zero.
Its coefficient of order $r$ is
\begin{equation}\label{eq:complexgross}
 [s^r]L_F^{(p)}(\chi,s)
 = L_F(\chi,0)\prod_{\mathfrak v\mid p,\ \mathfrak v\notin\mathcal R}
       (1-\chi(\mathfrak v))
   \prod_{\mathfrak v\in\mathcal R}\log N\mathfrak v .
\end{equation}
Indeed, $L_F(\chi,0)$ is nonzero by total oddness, the
functional equation and nonvanishing of the dual Hecke
$L$-value at $1$. The factors at $\mathcal R$ have simple
zeros with the indicated logarithmic slopes. This is the
complex Euler valuation determinant in the odd character
situation. Since $\log N\mathfrak v=f_\mathfrak v\log p$,
it singles out the same residue-degree weights as
\eqref{eq:normalizedmatrix}.

Let $\omega$ be the Teichm\"uller character of the rational
cyclotomic character restricted to $G_F$. The
Deligne--Ribet function in the cyclotomic variable,
$L_p(F,\chi\omega,s)$, uses the interpolation
$L_p(F,\chi\omega,n)=L_F^{(p)}(\chi\omega^n,n)$ for
integers $n\le0$ in the convention of \cite[eq.~(3)]{DKV}.
It is analytic near $0$. The algebraic value $L_F(\chi,0)$
is viewed $p$-adically using the same embedding as the
character. Define
\begin{equation}\label{eq:normcoeff}
 A(F,\chi)=
 \frac{[s^r]\,L_p(F,\chi\omega,s)}
 {L_F(\chi,0)
  \displaystyle\prod_{\mathfrak v\mid p,\,
                       \mathfrak v\notin\mathcal R}
  (1-\chi(\mathfrak v))
  \displaystyle\prod_{\mathfrak v\in\mathcal R}f_\mathfrak v}.
\end{equation}
All denominators are nonzero. Writing $[s^r]$ makes
explicit the division by $r!$ if derivatives are used.

\begin{theorem}\label{thm:analytic}
One has $\ord_{s=0}L_p(F,\chi\omega,s)\ge r(F,\chi)$ and
\begin{equation}\label{eq:factorcoeff}
 A(F,\chi)=\det B(F,\chi)
  =\prod_{\substack{\rho\text{ odd}\\d_\rho>0}}
     \det(C_\rho)^{m_{F,\chi}(\rho)} .
\end{equation}
If $d_\rho\le 1$ for every $\rho$ with $m_{F,\chi}(\rho)>0$,
then the order of vanishing is exactly $r(F,\chi)$.
\end{theorem}

\begin{proof}
The lower bound is the Gross--Stark vanishing bound, proved
cohomologically by Spiess and by Charollois and Dasgupta
\cite{Spiess,CD}. By Dasgupta--Kakde--Ventullo
\cite[Th.~1 and eq.~(8)]{DKV}, the quotient in
\eqref{eq:normcoeff} before division by the product of
residue degrees is $\det(L_\chi O_\chi^{-1})$.
This gives the first determinant in \eqref{eq:factorcoeff}.
The second equality is \eqref{eq:blockpoly}. When all the
contributing positive dimensions are one, each
$\det(C_\rho)$ is a nonzero single entry by
Proposition~\ref{prop:independent-end}. The normalized
Taylor coefficient is then nonzero, as required.
\end{proof}

\begin{remark}
The analytic formula uses the coefficient of the expected
order even when it vanishes. The missing nonvanishing
for a general higher-dimensional block is the
Gross--Kuz'min problem in the corresponding component.
Gross proves rank-one nonvanishing \cite[Prop.~2.13]{Gross}.
For a sum of one-dimensional multiplicity spaces the conclusion
follows as well from the rank bounds of
\cite[Thms.~4.3, 4.6]{Maksoud} and Corollary~\ref{cor:liftbound}. Formula
\eqref{eq:factorcoeff} gives the individual factors and
the exponent of each in any family of character projections.
\end{remark}

\begin{corollary}\label{cor:rankone}
Suppose $r(F_i,\chi_i)=1$ for each pair of a finite family.
There is a unique odd $\rho_i$ such that
$m_{F_i,\chi_i}(\rho_i)=d_{\rho_i}=1$. Then the
$A(F_i,\chi_i)$ are nonzero. Their affine linear relations
over $\Qbar$ are given by
\begin{equation}\label{eq:rankone-affine}
 b+\sum_i a_i A(F_i,\chi_i)=0
 \quad\Longleftrightarrow\quad
 b=0,\qquad \sum_{i:\rho_i=\rho}a_i=0\quad\text{for all }\rho .
\end{equation}
In particular each such coefficient is transcendental.
The ratio of two is $1$ if their labels $\rho_i$ are equal,
and is transcendental if their labels are different.
\end{corollary}

\begin{proof}
The dimension formula for $r$ implies the uniqueness.
By \eqref{eq:factorcoeff} the normalized coefficient is
the single entry of $C_{\rho_i}$. The assertion follows
from \eqref{eq:evaluation}. If a ratio of distinct entries
were algebraic, it would produce a nonzero linear relation.
\end{proof}

\begin{corollary}\label{cor:coefficientfield}
Let $\vartheta_\rho$ denote the character of $V_\rho$.
For every odd $\rho$ with $d_\rho>0$,
\[
 \Q_p(\Tr C_\rho)=\Q_p(\vartheta_\rho(g):g\in G).
\]
In particular an order-one coefficient $A(F,\chi)$
generates over $\Q_p$ the character field of its
unique contributing representation $\rho$.
\end{corollary}

\begin{proof}
Work in a finite Galois extension of $\Q_p$ containing
the coefficients needed to split the group algebra in
\eqref{eq:decomp}. The equivariant endomorphism $R_N$ is
defined over $\Q_p$. A Galois element $\sigma$ therefore
sends $\Tr C_\rho$ to $\Tr C_{\sigma\rho}$, where the
character of $\sigma\rho$ is $\sigma\vartheta_\rho$.
The integer $d_\rho$ is invariant under this action.
For distinct labels with positive $d_\rho$, the traces
are distinct by Proposition~\ref{prop:independent-end}
because they evaluate distinct coordinate functionals.
Thus the stabilizer of $\Tr C_\rho$ is exactly the
stabilizer of $\vartheta_\rho$. Galois correspondence
gives the field identity. When $r=1$, the distinguished
block has size one and its trace is $A(F,\chi)$ by
\eqref{eq:factorcoeff}.
\end{proof}

\subsection{Polynomial relations under weak Schanuel}

The weak $p$-adic Schanuel conjecture in the form of
\cite[Conj.~3.10]{CM} (see \cite[\S2]{Maksoud}) asserts that any finite
family of logarithms of nonzero algebraic numbers that is
linearly independent over $\Q$ is algebraically independent
over $\Q$. Algebraic independence over $\Q$ here implies
algebraic independence over $\Qbar$.

\begin{theorem}\label{thm:ideal}
For a finite family $(F_i,\chi_i)_{i=1}^n$, define
$m_i=(m_{F_i,\chi_i}(\rho))_{\rho\ {\rm odd},\,d_\rho>0}$.
Every binomial
\begin{equation}\label{eq:binomial}
 X^a-X^b,\qquad
 a,b\in\Z_{\ge0}^n,\qquad \sum_i a_i m_i=\sum_i b_i m_i ,
\end{equation}
vanishes at $(A(F_i,\chi_i))_{i=1}^n$ unconditionally.
Under weak $p$-adic Schanuel these binomials generate the
entire ideal of relations over $\Qbar$. The transcendence
degree of the field generated by the coefficients over
$\Qbar$ is then the rational rank of the vectors $m_i$.
\end{theorem}

\begin{proof}
The unconditional statement follows by multiplying the
determinant identities \eqref{eq:factorcoeff}. For the
conditional direction, the compressed algebra $\mathcal A$
in Proposition~\ref{prop:independent-end} is defined over
$\Q$. Evaluate the logarithm on the $O_N^{-1}$ image of
a rational basis of $\mathcal A$. Clearing the rational
denominators of those unit vectors gives logarithms of
actual algebraic numbers. They are rationally independent
by Lemma~\ref{lem:log}. The hypothesis makes the
logarithms algebraically independent. The change from that
basis of evaluations to the full set of block entries of
\eqref{eq:decomp} is invertible over $\Qbar$. Hence the
block entries are algebraically independent.

For matrices of independent variables indexed by $\rho$,
the determinants of the separate blocks are algebraically
independent. To see this, specialize each block to the
diagonal matrix with first entry $z_\rho$ and remaining
entries $1$. The determinants specialize to independent
variables $z_\rho$. Consequently the only relations among
\eqref{eq:factorcoeff} come from the homomorphism
$\Qbar[X_1,\ldots,X_n]\to\Qbar[z_\rho]$ sending
$X_i$ to $z^{m_i}$. Monomials in the $z_\rho$ are linearly
independent. Grouping a polynomial by the image exponent
shows that its kernel is generated by \eqref{eq:binomial}.
The fraction field of the monomial algebra has
transcendence degree equal to the rank of the lattice
spanned by the $m_i$, by choosing a lattice basis of its
group of Laurent monomials.
\end{proof}

\section{Polynomial relations among the cyclotomic \texorpdfstring{$L$}{L}-functions}
\label{sec:germs}

The constant $A(F,\chi)$ involves only the Artin constituents having
a fixed vector under the decomposition group at $p$. Which
constituents can be recovered from the analytic function before
evaluation at zero? In this section the answer is the full induced
character. The argument uses interpolation at negative integers and
the complex functional equation. It applies independently of any
conjecture on logarithms.

Continue to fix $p$, $N$ and $G$ as in Section~\ref{sec:galois}.
A germ at zero is a convergent function on some $p$-adic neighborhood
of zero, with two functions identified when they agree on a smaller
neighborhood. Let $\mathcal M_p$ be the field of meromorphic
$\Cp$-valued germs and let
\[
 \mathcal K_p=\Cp\bigl(s,\,\langle u\rangle^s:
       u\in\Q_{>0}^{\times},\ \ord_p(u)=0\bigr)\subset\mathcal M_p .
\]
Here $\langle u\rangle=u/\omega(u)$ and
$\langle u\rangle^s=\exp_p(s\log_p u)$. This coefficient field
contains the usual Iwasawa parameter $(1+p)^s-1$ and the
cyclotomic Euler factors at primes away from $p$.

For a finite family of pairs write
\[
 V_i=\operatorname{Ind}_{\Gal(N/F_i)}^G\chi_i,\qquad
 \mathcal L_i(s)=L_p(F_i,\chi_i\omega,s),\qquad r_i=r(F_i,\chi_i).
\]
We use $\chi_i$ in $V_i$ because Artin induction acts on the character
of the complex $L$-function. The multiplicities $m_{F_i,\chi_i}(\rho)$
in \eqref{eq:decomp} concern $V_i^\vee$. Equality of two sums of
these representations is equivalent to equality of their duals.
Let $c_i$ be the entire denominator of \eqref{eq:normcoeff} and put
\[
 \mathcal H_i(s)=\frac{\mathcal L_i(s)}{c_i s^{r_i}}.
\]
These are analytic germs with $\mathcal H_i(0)=A(F_i,\chi_i)$,
including when that coefficient is zero.

\begin{theorem}[Artin relations for cyclotomic germs]\label{thm:germrelations}
The ideal of polynomial relations over $\mathcal K_p$ among the
$\mathcal L_i$ is generated by
\begin{equation}\label{eq:germbinomials}
 X^a-X^b,\qquad a,b\in\Z_{\ge0}^n,\qquad
 \bigoplus_i V_i^{\oplus a_i}\simeq\bigoplus_i V_i^{\oplus b_i}.
\end{equation}
The same binomials generate the ideal for the $\mathcal H_i$.
For either family the transcendence degree over $\mathcal K_p$ is
the rational rank of the induced characters in the character ring.
In particular a family with linearly independent induced characters
has algebraically independent cyclotomic $L$-functions.
\end{theorem}

Artin formalism in the $p$-adic setting is classical
\cite{GreenbergArtin}. For complex functions, Kaczorowski, Molteni
and Perelli prove linear independence in classes of Euler products
\cite{KMP}, and Cimpoea\c{s} and Nicolae prove algebraic independence
for irreducible Artin characters over meromorphic functions of order
less than one \cite[Cor.~9]{CN}. L\"u permits a growth bound along
a sequence of radii \cite[Cor.~2]{Lu}. For complex Hasse--Weil
functions of abelian varieties, Yoshikawa obtains a multiplicity
criterion in terms of simple isogeny factors \cite[Th.~1.3]{Yoshikawa}.
Angl\`es and Ranieri prove linear independence for
reductions of Dirichlet Iwasawa series over a field of
pseudo-rational functions in characteristic $p$ \cite[Th.~1]{AR}.
Theorem~\ref{thm:germrelations} concerns the characteristic-zero
interpolating germs themselves. We give the proof because the
passage between the two absolute values and the extra functions in
$\mathcal K_p$ both matter here.

\begin{lemma}[Values at interpolation integers]\label{lem:archsamples}
Fix the complex embedding of $\Qbar$ used for the Artin characters.
Let $V$ be a direct sum of the $V_i$, allowing the zero representation. Write $d_V=\dim V$, $Q_V$ for its positive Artin
conductor over $\Q$, and $h_V=\dim V^{I_p}$. Let
$P_{V,p}(X)=\det(1-X\operatorname{Frob}_p\mid V^{I_p})$ in the
Artin convention for $L(V,s)$. There is a constant $\xi_V\ne0$
such that, for every positive $n$ divisible by $4$,
\begin{align}
 L^{(p)}(V,-n)
  &=\xi_V\left(\frac{n!}{2^n\pi^{n+1}}\right)^{d_V}
          (Q_Vp^{h_V})^n D_V(n),\label{eq:sampleformula}\\
 D_V(z)&=P_{V^\vee,p}(p^{-z}) L(V^\vee,1+z)
        =\sum_{m\ge1} b_V(m)m^{-z}\quad(\Re z>0).\nonumber
\end{align}
The series is absolutely convergent with $b_V(1)=1$. For $m$ a
squarefree product of primes unramified in $N$ and different from
$p$, its coefficient is
\begin{equation}\label{eq:squarefree}
 m b_V(m)=\prod_{q\mid m}\operatorname{Tr}
                     (V^\vee(\operatorname{Frob}_q)).
\end{equation}
\end{lemma}

\begin{proof}
The completed function over $\Q$ has archimedean factor
$\Gamma_\R(s+1)^{d_V}$, because complex conjugation acts as $-1$,
where $\Gamma_\R(z)=\pi^{-z/2}\Gamma(z/2)$. Use the functional
equation for $Q_V^{s/2}\Gamma_\R(s+1)^{d_V}L(V,s)$ with
root number $w_V$ \cite[Ch.~VII, \S12]{Neukirch}.
At $s=-n$ the ratio $\Gamma_\R(n+2)/\Gamma_\R(1-n)$ is
$n!/(2^n\pi^{n+1})$ by reflection and duplication, with positive
sign for $4\mid n$. To remove the factor at $p$, multiply by
$P_{V,p}(p^n)$. Finite inertia gives
$(V^\vee)^{I_p}=(V^{I_p})^\vee$ canonically by averaging, so
\[
 P_{V,p}(p^n)=
 \det(-\operatorname{Frob}_p\mid V^{I_p})\,p^{nh_V}
                     P_{V^\vee,p}(p^{-n}).
\]
This proves \eqref{eq:sampleformula} with
$\xi_V=w_VQ_V^{1/2}\det(-\operatorname{Frob}_p\mid V^{I_p})$.
Absolute convergence follows from the complex Euler product in
$\Re(1+z)>1$. Its coefficients are multiplicative, and the Euler
polynomial multiplying it has support on powers of $p$.
This gives \eqref{eq:squarefree}, including the empty product.
\end{proof}

\begin{proof}[Proof of Theorem~\ref{thm:germrelations}]
For a direct sum $V=\bigoplus_iV_i^{\oplus a_i}$ put
$\mathcal L_V=\prod_i\mathcal L_i^{a_i}$. At the integers
\[
                s=-n,\qquad n=4(p-1)p^j\quad(j\ge1),
\]
interpolation gives $\mathcal L_V(-n)=\iota L^{(p)}(V,-n)$,
using Artin induction and $\omega^n=1$. Here $\iota$ is our fixed
embedding of algebraic numbers into $\Cp$. These $L$-values are
algebraic by Siegel's theorem, in the interpolation formulation
of \cite[eq.~(3) and the following paragraph]{DKV}. They are nonzero by
\eqref{eq:sampleformula} and the Euler product at $1+n$.
Since $-n\to0$ $p$-adically, the identity principle proves that
$\mathcal L_V$ depends only on $V$. This gives every relation
\eqref{eq:germbinomials}.

We prove that the germs $\mathcal L_V$ for pairwise nonisomorphic
such sums $V$ are linearly independent over $\mathcal K_p$.
After clearing denominators, a relation would give a finite sum
\begin{equation}\label{eq:clearedgerms}
 \sum_{(V,u)} P_{V,u}(s)\langle u\rangle^s\mathcal L_V(s)=0 ,
\end{equation}
where the pairs are distinct, $u$ is a positive rational $p$-unit,
and the $P_{V,u}$ are polynomials, some nonzero. Products of the
finite exponential factors have simply been combined by
multiplying their rational bases. At the chosen integers the
exponential factor is $u^{-n}$. For any fixed bound on the degrees
of the polynomials, the equations for their coefficients at all
these integers have entries in $\iota(\Qbar)$. The span of the rows
has a finite basis consisting of rows of the same matrix. Its rank
over $\iota(\Qbar)$ equals its rank over $\Cp$, since rank is
detected by finite minors. A nonzero solution would therefore give
a nonzero solution over $\Qbar$ after transporting the equations
to a complex embedding. Discard pairs with zero polynomial in this
solution. We use only the resulting complex identities at the chosen
integers.

Let $d$ be the largest $d_V$ among the remaining pairs. Divide
by $(n!/(2^n\pi^{n+1}))^d$ and put $A_{V,u}=Q_Vp^{h_V}/u$.
The terms with $d_V<d$ are $o(c^n)$ for every $c>0$ as
$j\to\infty$: the factorial dominates every fixed exponential and
$D_V(n)$ is bounded. For $d_V=d$, expand $D_V(n)$ as in
\eqref{eq:sampleformula}. The positive frequencies are
$x=A_{V,u}/m$, $m\ge1$. There are only finitely many above any
positive bound. We claim that for each such $x$ one has the
polynomial identity
\begin{equation}\label{eq:frequency}
 \sum_{\substack{d_V=d,\ m\ge1\\ A_{V,u}/m=x}}
      \xi_V b_V(m)P_{V,u}(X)=0 .
\end{equation}
If a nonzero polynomial occurred, choose the largest corresponding
frequency $x$. Its contribution, divided by $x^n$, would be that
polynomial at $-n$. All smaller frequencies contribute
$O(n^k(x'/x)^n)$ for some $x'<x$ and the chosen degree bound $k$.
Indeed, for each of the finitely many Dirichlet series, sum the
tail by reserving one fixed positive exponent of absolute
convergence and bound all remaining powers by $(x')^{n-O(1)}$.
The lower-dimensional terms tend to zero after the same division.
A nonzero polynomial at $-n$ cannot tend to zero. This proves
\eqref{eq:frequency}.

Choose $A_*$ maximal among the $A_{V,u}$ with $d_V=d$. Take
$m$ squarefree, supported on unramified primes away from $p$
and from all denominators of the reduced fractions $A_{V,u}/A_*$.
At $x=A_*/m$ only the pairs with $A_{V,u}=A_*$ contribute:
a smaller ratio has denominator greater than one and its product
with $m$ is not an integer. For a fixed $V$ at this maximal value
there is at most one $u$. Equations \eqref{eq:frequency} and
\eqref{eq:squarefree} therefore give
\[
 \sum_{(V,u):\,d_V=d,\ A_{V,u}=A_*}
 \xi_V P_{V,u}(X)
       \prod_{q\mid m}\operatorname{Tr}(V^\vee(\operatorname{Frob}_q))=0 .
\]
Chebotarev's density theorem
\cite[Ch.~VII, Th.~13.4]{Neukirch} allows any prescribed finite list of
conjugacy classes, including repetitions, to be represented by
distinct primes avoiding this finite excluded set. The last equality
thus holds for every monomial in the character coordinates on the
conjugacy classes of $G$. The character tuples for these $V^\vee$ are
distinct, since finite group representations in characteristic zero
are determined by their characters. Polynomials separate finitely many
distinct tuples. Fix a tuple and, for each other tuple, choose a
coordinate on which the two differ. A normalized linear polynomial
in that coordinate is one at the chosen tuple and zero at the other.
The product of these linear polynomials isolates the chosen tuple. Substituting this polynomial gives
$\xi_VP_{V,u}(X)=0$ for the chosen pair, a contradiction.

Group any polynomial in the $\mathcal L_i$ by the isomorphism
class of the associated direct sum. The independence just proved
forces the coefficient sum of each group to vanish, so the kernel
is generated by differences of monomials in the same group.
For the $\mathcal H_i$, the product of the normalizing factors
$c_i s^{r_i}$ depends only on the direct sum. Indeed, by
\eqref{eq:complexgross} the order of the complex $p$-depleted
function is $\sum_i a_i r_i$, and its leading coefficient is
$(\log p)^{\sum_i a_i r_i}\prod_i c_i^{a_i}$.
Thus each group is scaled by a single nonzero element of
$\mathcal K_p$. Finally, the fraction field of the resulting
monomial algebra has transcendence degree the rational rank of its
exponent vectors, by the lattice-basis argument in
Theorem~\ref{thm:ideal}.
\end{proof}

\begin{corollary}\label{cor:loss}
Write $v_i=(m_{F_i,\chi_i}(\rho))_{\rho\ {\rm odd}}$ and let
$\pi_D$ discard the coordinates with $d_\rho=0$. The germs have
transcendence degree $\rank_\Q\langle v_i\rangle$ over $\mathcal K_p$.
Under weak $p$-adic Schanuel the values $\mathcal H_i(0)$ have
transcendence degree $\rank_\Q\langle\pi_Dv_i\rangle$ over $\Qbar$.
The difference of these two ranks is
$\dim_\Q(\langle v_i\rangle_\Q\cap\ker\pi_D)$.
\end{corollary}

\begin{proof}
Apply Theorem~\ref{thm:germrelations} to the germs and
Theorem~\ref{thm:ideal} to the constants. The last assertion is
rank-nullity on the indicated rational character space.
\end{proof}

For instance, if $\psi$ is an odd Dirichlet character with
$\psi(p)\ne1$, its normalized germ has value $1$ at zero by
interpolation. The germ is transcendental over $\mathcal K_p$
by Theorem~\ref{thm:germrelations}. The next section gives an
example in which this phenomenon accounts for the first failure
of an identity in the Taylor expansions.

\subsection{Irreducible \texorpdfstring{$p$}{p}-adic Artin functions}

For a totally odd representation $V$ of $G$ let $\mathcal L_V$ denote
the meromorphic $p$-adic Artin germ for $V\otimes\omega$ in the
interpolation convention above. Its construction for a nonmonomial
$V$ uses Brauer induction \cite{GreenbergArtin}. Put
\[
 r(V)=\dim V^D,\qquad
 c(V)=\frac{[s^{r(V)}]L^{(p)}(V,s)}{(\log p)^{r(V)}} .
\]
The complex definition of $c(V)$ gives a nonzero algebraic number,
as the proof below shows. We use its image under $\iota$ in a
$p$-adic formula.

\begin{corollary}[Artin germs and irreducible regulators]\label{cor:artingerms}
The germs $\mathcal L_V$ for the distinct odd irreducible
representations $V$ of $G$ are algebraically independent over
$\mathcal K_p$. In particular, distinct irreducible totally even
Artin representations of the absolute Galois group $G_\Q$ have algebraically independent
meromorphic $p$-adic $L$-germs in this convention.

Suppose in addition that each nonempty block $C_\rho$ for $N$ is
invertible. For every totally odd $V$, the order of $\mathcal L_V$
at zero is $r(V)$ and
\begin{equation}\label{eq:allartinleading}
 \left.\frac{\mathcal L_V(s)}{c(V)s^{r(V)}}\right|_{s=0}
   =\prod_{\rho:\,d_\rho>0}(\det C_\rho)^{[V^\vee:\rho]} .
\end{equation}
Here $[V^\vee:\rho]$ is the multiplicity of $\rho$ in $V^\vee$.
The invertibility hypothesis holds if every odd $d_\rho\le1$.
It holds for any $N$ under weak $p$-adic Schanuel.
\end{corollary}

\begin{proof}
Brauer's induction theorem writes $V$ as an integral virtual
sum of characters induced from linear characters $\lambda$ of
subgroups $J\subseteq G$ \cite[Ch.~VII, \S10]{Neukirch}.
Take the odd part of this identity. If $j\in J$, retain the
summand precisely when $\lambda(j)=-1$. If $j\notin J$, then
$J\langle j\rangle=J\times\langle j\rangle$, and induction
from $J$ to this subgroup is the sum of the two extensions of
$\lambda$ with opposite values at $j$. Its odd part is induced
from the extension taking $j$ to $-1$. Thus all terms in the
resulting virtual sum are induced from odd linear characters over
totally real fixed fields. They are exactly of the kind used in
Theorem~\ref{thm:germrelations}.

Taking the corresponding product and quotient of the monomial
germs constructs $\mathcal L_V$, independently of the chosen
decomposition by \eqref{eq:germbinomials}. To test a polynomial
relation for the irreducibles, clear the finitely many monomial
denominators. Distinct monomials in distinct irreducibles have
distinct characters even after adding the common denominator
character. The linear independence in the proof of
Theorem~\ref{thm:germrelations} therefore forces every coefficient
to vanish. For any finite family of totally even representations
over $\Q$, take the compositum of their totally real splitting
fields with $\Q(\mu_p)$. Twisting the family by $\omega^{-1}$
gives odd representations of this CM extension. Twisting back
recovers the asserted $p$-adic germs.

Write $[V]=\sum_i n_i[V_i]$ for the odd Brauer decomposition.
Local invariants give $r(V)=\sum_i n_i r_i$ by Frobenius
reciprocity. Complex Artin formalism and \eqref{eq:complexgross}
give $c(V)=\prod_i c_i^{n_i}\in\Qbar^\times$. If all the
blocks are invertible, Theorem~\ref{thm:analytic} shows that every
$\mathcal H_i$ is a unit at zero. The quotient
$\prod_i\mathcal H_i^{n_i}$ is therefore regular and nonzero at
zero, where its value is the right side of
\eqref{eq:allartinleading} by additivity of multiplicities.
The two sufficient conditions follow from
Proposition~\ref{prop:independent-end} in dimension one and the
algebraic independence of the block entries in the proof of
Theorem~\ref{thm:ideal} in arbitrary dimension.
\end{proof}

\section{Nonabelian examples with explicit coefficient identities}
\label{sec:examples}

Suppose $N=L I$, where $L/\Q$ is totally real Galois with
group $S_3$ and $I/\Q$ is imaginary quadratic. Suppose $p$
is unramified in $N$, splits in $I$, and has Frobenius a
transposition in $L$. Then $G=S_3\times C_2$ and $D$
is generated by $(\tau,1)$ for a transposition $\tau$.
Write $\epsilon$ for the odd character of the $C_2$ factor.
The odd irreducibles are $1\otimes\epsilon$,
$\mathrm{sgn}\otimes\epsilon$, and
$\mathrm{std}\otimes\epsilon$. Their $D$-fixed dimensions
are $1,0,1$, respectively, since the standard representation
of $S_3$ has eigenvalues $1,-1$ on a transposition.
Write $\alpha,\beta$ for the entries of the corresponding
one-dimensional blocks for $1\otimes\epsilon$ and
$\mathrm{std}\otimes\epsilon$.

Here $\mathrm{std}$ denotes the two-dimensional sum-zero summand
of the permutation representation on three letters.
As a character of $G$ over $\Q$, $\epsilon$ gives
$A(\Q,\epsilon)=\alpha$. Let $F_2$ be the real quadratic
subfield of $L$. Then $\Gal(N/F_2)=C_3\times C_2$.
For either nontrivial character $\theta$ of $C_3$, the
induction of $(\theta\epsilon)^{-1}$ to $G$ is
$\mathrm{std}\otimes\epsilon$. Hence
$A(F_2,\theta\epsilon)=\beta$ and the expected order is one.
The two choices of $\theta$ give the same normalized coefficient.
By Corollary~\ref{cor:rankone}, $1,\alpha,\beta$ are
$\Qbar$-linearly independent.

For a non-Galois cubic subfield $F_3$ of $L$, let
$\epsilon_{F_3}$ be the character of $F_3I/F_3$. Its
induction contains each of $1\otimes\epsilon$ and
$\mathrm{std}\otimes\epsilon$ once. For $\epsilon_L$ of
$LI/L$, the multiplicities of these two irreducibles
are $1$ and $2$. This can be checked by restriction to
the appropriate subgroup. For $F_3$ the subgroup in
$S_3$ is a transposition group, and for $L$ it is
trivial. Theorem~\ref{thm:analytic} now gives the
following identities for the normalized Taylor coefficients.
\begin{equation}\label{eq:s3coeff}
 A(F_3,\epsilon_{F_3})=\alpha\beta,\qquad
 A(L,\epsilon_L)=\alpha\beta^2,\qquad
 A(L,\epsilon_L)A(\Q,\epsilon)=A(F_3,\epsilon_{F_3})^2 .
\end{equation}
The orders of vanishing for the cubic and normal fields
are exactly two and three. In the unscaled convention
for the Gross determinant of \cite[eq.~(8)]{DKV}, the
residue-degree product is $2$ for $F_3$ and $2^3$ for
$L$. The last identity of \eqref{eq:s3coeff} therefore
acquires a factor of $2$ on its right in that convention.

\begin{example}\label{ex:concrete}
The splitting field $L$ of $x^3-4x+1$ is totally real and
has group $S_3$. Indeed the polynomial is irreducible by
the rational root test, and its discriminant is $229>0$,
which is not a square. At $p=7$ the polynomial factors as
$(x-4)(x^2+4x+5)$ modulo $7$, with the quadratic
irreducible. Thus Frobenius is a transposition. The prime
$7$ splits in $I=\Q(\sqrt{-3})$. For the embedding through
the prime generated by $\pi=2+(1+\sqrt{-3})/2$, of norm
$7$, one has
\[
 \alpha=-\log_7(\iota(\pi/\overline{\pi})).
\]
The ratio has valuation one at the selected prime, so
this equality follows directly from the one-dimensional
Gross matrix. The other scalar $\beta$ is the normalized
coefficient for either of the order-six characters of
$\Gal(N/F_2)$ used above.
\end{example}

\subsection{The next coefficient at \texorpdfstring{$p=7$}{p = 7}}
\label{subsec:nextcoefficient}

Keep the fields of Example~\ref{ex:concrete}. Put
$\delta=\mathrm{sgn}\otimes\epsilon$, the quadratic Dirichlet
character of $\Q(\sqrt{-687})$. Use the subscripts $L,3,1,\delta$
on $\mathcal H$ for $(L,\epsilon_L)$, $(F_3,\epsilon_{F_3})$,
$(\Q,\epsilon)$, and $(\Q,\delta)$ respectively. Their expected
orders are $3,2,1,0$. The character $\delta$ has conductor $687$,
since $-687=-3\cdot229$ is a fundamental discriminant. Its value
at $7$ is $-1$.

\begin{proposition}\label{prop:next}
The normalized germs satisfy
\[
 \frac{\mathcal H_L(s)\mathcal H_1(s)}{\mathcal H_3(s)^2}
       =\mathcal H_\delta(s)
       =\frac{L_p(\Q,\delta\omega,s)}{2L(\delta,0)} .
\]
The right side equals $1$ at zero and has derivative congruent to
$42$ modulo $49$ in $\Z_7$. In particular the last identity in
\eqref{eq:s3coeff} holds for the constant coefficients and fails
for the linear coefficients of these normalized Taylor series.
\end{proposition}

\begin{proof}
As characters of $S_3$, the regular representation plus the trivial
representation is twice the permutation representation on the three
cosets of a transposition subgroup plus the sign representation.
Tensor with $\epsilon$ and apply Theorem~\ref{thm:germrelations}.
The normalization $2L(\delta,0)$ follows from $\delta(7)=-1$.
The nonvanishing in \eqref{eq:s3coeff} permits division by
$\mathcal H_3$ near zero.

Here is an exact computation of the slope. Set $f=687$,
$S_j=\sum_{a=1}^f\delta(a)a^j$, and $F(s)=L_p(\Q,\delta\omega,s)$.
Quadratic reciprocity gives
$\delta(a)=(\frac a3)(\frac a{229})$, with each Legendre symbol
zero on multiples of its prime. For $k\ge1$ the generalized
Bernoulli formula \cite[Ch.~4]{Washington} reads
\[
 L(\delta,1-k)=-\frac{f^{k-1}}{k}\sum_{a=1}^f\delta(a)B_k(a/f).
\]
For $k=1,7$ this gives
\begin{align*}
 L(\delta,0)&=-S_1/f=12,\\
 L(\delta,-6)&=-\frac{6S_7-21fS_6+21f^2S_5-7f^4S_3+f^6S_1}{42f}
                  =-10344805682225424 .
\end{align*}
These two sums can be checked using the nonzero squares modulo
$3$ and $229$. The first sum gives $S_1=-8244$.
Interpolation gives $F(0)=24$ and
$F(-6)=(1+7^6)L(\delta,-6)\equiv3\pmod{49}$.

The even Dirichlet character $\delta\omega$ is nontrivial on the
prime-to-$7$ conductor part and is of the first kind. Its integral
Iwasawa series gives $F(s)=H(8^s-1)$ for some $H\in\Z_7[[T]]$
\cite[Th.~7.10]{Washington}. The values of $\delta\omega$ lie in
$\Z_7$. A change of topological generator of $1+7\Z_7$ acts by
$1+T\mapsto(1+T)^u$ for $u\in\Z_7^\times$. Translating by an
integral center or reversing the cyclotomic variable preserves
integrality. For $s\in\Z_7$
we have $8^s-1\equiv 7s\pmod{49}$ and
$\log_7(8)\equiv7\pmod{49}$. It follows that
$F(s)-F(0)\equiv sF'(0)\pmod{49}$. Taking $s=-6$ gives
$F'(0)\equiv28\pmod{49}$. Division by $24$ gives
$\mathcal H_\delta'(0)\equiv42\pmod{49}$.
\end{proof}

Thus a constituent on which $D$ has no fixed vector contributes
to the linear coefficient even though its normalized constant is
one. Theorem~\ref{thm:germrelations} gives a second description:
$\mathcal H_\delta$ is transcendental over $\mathcal K_7$.

\subsection{Fields of definition and growth in towers}

For a coefficient whose field of definition shrinks,
take the same cubic splitting field $L$ but use $p=29$
and $I=\Q(i)$. The discriminant $229$ is congruent to
$-3$ modulo $29$, a quadratic nonresidue because
$29\equiv-1\bmod 3$. Therefore Frobenius in $L$ is a
transposition. The prime splits in $I$ since
$29\equiv1\bmod4$. For $\chi=\theta\epsilon$ over $F_2$,
the character values generate $\Q_{29}(\mu_3)$ over
$\Q_{29}$. The order-one coefficient lies in $\Q_{29}$
by Corollary~\ref{cor:coefficientfield}, since
$\mathrm{std}\otimes\epsilon$ has rational character.

The same computation gives unbounded exact orders at $p=7$.
Keep $L$ and $I$ as in
Example~\ref{ex:concrete}. Choose distinct primes
$q_1,\ldots,q_m$ congruent to $1$ modulo $28$ and different
from $229$. Their real quadratic fields have a compositum
$K_m$ with group $C_2^m$ and with $7$ split completely.
The only quadratic subfield in $L$ is $\Q(\sqrt{229})$,
whereas every quadratic subfield of $K_m$ has a square
class supported on the $q_i$. Since every abelian
subextension of the $S_3$ extension has degree at most
two, $K_m$ is linearly disjoint from $L$. The compositum
$LK_m$ is totally real and has trivial intersection with
$I$. Therefore $N_m=LK_mI$ has group $S_3\times C_2^m
\times C_2$, with $D=\langle(\tau,1,1)\rangle$ at $7$.
There are arbitrarily many choices of the $q_i$ by
Dirichlet's theorem on primes in progressions.

For each character $\eta$ of $C_2^m$, the positive
$D$-fixed dimensions in the odd sector attached to
$\eta$ are both one. They correspond to
$1\otimes\eta\otimes\epsilon$ and
$\mathrm{std}\otimes\eta\otimes\epsilon$.
Let $\alpha_\eta$ be the normalized coefficient over
$\Q$ for the former character. Let $\beta_\eta$ be the
normalized coefficient over $F_2$ for
$\theta\otimes\eta\otimes\epsilon$ on
$\Gal(N_m/F_2)=C_3\times C_2^m\times C_2$.
These are all coefficients of order one. Their
distinct irreducible labels imply that the entire
family is linearly independent over $\Qbar$ together
with $1$. For the quadratic character of $N_m/N_m^+$,
the induction has each odd irreducible with multiplicity
equal to its dimension. It follows that its normalized
coefficient is
\[
  \prod_{\eta\in\widehat{C_2^m}}\alpha_\eta\beta_\eta^2
\]
and its order of vanishing is $3\cdot2^m$.

\section{Adding an abelian part}
\label{sec:limits}

There is a direct extension of the complex theorem that uses
no new special-value input for the torus. Let $A/\Q$ be an
abelian variety and let $M'$ be an extension
\begin{equation}\label{eq:abelianext}
 0\longrightarrow M_S(K/k)\longrightarrow M'
 \longrightarrow [0\to A]\longrightarrow0
\end{equation}
with torus and lattice as indicated. Thus the semiabelian
part of $M'$ may be a nonsplit extension and the given
lattice map lands in its torus. Alternatively take
$G=T_{K/k}\times A$ and allow any abelian component of
the lattice map, keeping its toric component equal to
that of $M_S(K/k)$.

\begin{proposition}\label{prop:abelian}
In either case, the augmented real comparison of $M'$ is
an isomorphism. Its determinant is
\[
 c_S^+(M')=c_S^+(M_S(K/k))\,c^+([0\to A]).
\]
Moreover
\[
 L_{\mathrm{inv}}(M',s)=L_{\mathrm{inv}}(M_S(K/k),s)
                  L(A,s+1)
\]
with the usual cohomological $L$-function of $A$ on the right.
When $A$ is an elliptic curve and $L(A,1)\ne0$, the
leading coefficient for $M'$ represents $c_S^+(M')$ in
$\R^\times/\Q^\times$, and its order at zero is $d_S$.
\end{proposition}

\begin{proof}
For \eqref{eq:abelianext}, use exactness of the realizations,
the real plus functor, and the Lie sequence. The finite
valuation rows are those of $M_S(K/k)$. The two diagonal
blocks in bases from these exact sequences are the augmented
period map of $M_S(K/k)$ and the period map for $A$, both
isomorphisms. In the split semiabelian case the motive $[0\to A]$
is a subobject with quotient $M_S(K/k)$, giving the same
determinant equality with the two blocks in reverse order.
For both cases the toric valuation of each lattice vector
agrees with that of $M_S(K/k)$: in the first case by
restriction to $T_{K/k}$ and in the second because
$A(\Q_q)$ is compact. Formula \eqref{eq:globaleuler}
gives the identity of $L$-functions. For an elliptic curve
over $\Q$, modularity and the rationality of modular
symbols give $L(A,1)/c^+([0\to A])\in\Q$
\cite{MTT,Stevens,Wiles,TaylorWiles,BCDT}. The period in a rational differential and the N\'eron real period
differ by a nonzero rational factor. This proves the assertion for
the stated nonzero central value.
\end{proof}

For positive analytic rank of an elliptic curve, the
unconditional identity above isolates the additional
factor as the leading term of $L(A,s+1)$.
The BSD prediction for this factor uses a height
regulator. In particular, taking $K=k$ leaves
$[0\to A]$ with $L_{\mathrm{inv}}(M',s)=L(A,s+1)$.
A leading-term theorem for every 1-motive with its
conjectured arithmetic normalization would include
the corresponding BSD statement for elliptic curves.

There is a parallel restriction at $p$. For an elliptic
curve $E/\Q$ with good ordinary reduction at $p$, let
$u_p$ be the unit root of $X^2-a_pX+p$. In the cyclotomic
variable centered at the complex point $s=1$, ordinary
interpolation gives
\[
 L_p(E,0)=(1-u_p^{-1})^2\,L(E,1)/\Omega_E^+ ,
\]
where the N\'eron period and the $p$-adic function use
compatible rational normalizations \cite{MTT}.
Multiplying this identity by the relative $p$-adic zeta
quotient in \eqref{eq:padicS} evaluates the coefficient
of order $d_S$ of that product. When $L(E,1)\ne0$ and the
two totally real fields satisfy Leopoldt, the order is
exactly $d_S$. Indeed $u_p\ne1$ since
$p+1-a_p=\#E(\mathbf F_p)>0$.

For split multiplicative reduction the interpolation factor
produces an additional zero. The exceptional-zero theorem of
Greenberg and Stevens \cite[Th.~0.3, $p\ge5$]{GS}, with the
Coleman-map proof of Kobayashi \cite[Th.~4.1 and Cor.~4.2]{Kobayashi}
for odd $p$, gives in the same centered variable
\[
 \left.\frac{d}{ds}L_p(E,s)\right|_{s=0}
   = \frac{\log_p q_E}{\ord_p q_E}\frac{L(E,1)}{\Omega_E^+},
\]
where $q_E$ is the Tate parameter. The numerator is nonzero. By the $p$-adic Mahler--Manin theorem
\cite{BDGP}, algebraicity of $j(E)$ forces $q_E$ to be
transcendental over $\Q$. If $\log_p q_E=0$, the element
$q_E/p^{\ord_p q_E}$ of $\Z_p^\times$ would be a root of unity,
contradicting this transcendence. Consequently, if $L(E,1)\ne0$
and the hypotheses for exact order in Corollary~\ref{cor:padicS}
hold, the explicit product
$\mathcal Z_{K/k,S,p}(s)L_p(E,s)$ has order exactly $d_S+1$.
Its leading coefficient is the product of the coefficient in that
corollary and the displayed elliptic derivative. At positive rank, the
formulation of Mazur, Tate and Teitelbaum requires the
cyclotomic $p$-adic height regulator. The local height
pairings for 1-motives of Rivera Arredondo \cite{Rivera}
give a geometric framework for such additional terms.

The Tate parameter has the same two Kummer coordinates as the
characters in Theorem~\ref{thm:cuplift}. The elliptic convention uses
the positive logarithm, whereas \eqref{eq:LN} uses its negative.
The elliptic interpolation uses modular symbols. The odd Artin
interpolation uses the finite-order Hecke character. In Kobayashi's
calculation the local dual class annihilates $\kappa(q_E)$, and
this forces the quotient of logarithm and valuation
\cite[eqs.~(3)--(6)]{Kobayashi}. This is the one-column lifting
calculation in concrete elliptic form. Besser and de Shalit study
cohomological and analytic $\mathcal L$-invariants for varieties
uniformized by Drinfeld's symmetric domain \cite{BesserShalit}.

More generally, the $p$-adic special-value framework of
Perrin-Riou and the treatment of trivial zeros by Benois
depend on a local choice of regular submodule and on a
global regulator \cite{PerrinRiou,Benois}. The present
$p$-adic formulas concern two specified analytic functions.
In the totally real case the comparison is the finite
Kummer logarithm with Euler factors \eqref{eq:padicS}.
In the totally odd case it compares the valuation and
the norm logarithm of semistable Kummer classes, with
the interpolation in \eqref{eq:normcoeff}. A single
construction for arbitrary mixed 1-motives would have to
account for both of these situations and for the abelian
height factor. Even in the odd toric case, the determinant
of a block with $d_\rho\ge2$ presents a genuine
nonvanishing problem beyond linear independence of its
entries. Theorem~\ref{thm:ideal} states explicitly the
additional algebraic independence hypothesis that settles
all such block determinants and controls their relations.

\begin{thebibliography}{99}

\bibitem{AR}
B. Angl\`es and G. Ranieri, On the linear independence of $p$-adic
$L$-functions modulo $p$, \emph{Ann. Inst. Fourier (Grenoble)}
\textbf{60} (2010), 1831--1855.
\href{https://doi.org/10.5802/aif.2573}{DOI}.

\bibitem{Baker}
A. Baker, Linear forms in the logarithms of algebraic numbers. I,
\emph{Mathematika} \textbf{13} (1966), 204--216.

\bibitem{BDGP}
K. Barr\'e-Sirieix, G. Diaz, F. Gramain and G. Philibert,
Une preuve de la conjecture de Mahler--Manin,
\emph{Invent. Math.} \textbf{124} (1996), 1--9.
\href{https://doi.org/10.1007/s002220050044}{DOI}.

\bibitem{Beilinson}
A. A. Beilinson, Higher regulators and values of $L$-functions,
\emph{J. Soviet Math.} \textbf{30} (1985), 2036--2070.
\href{https://doi.org/10.1007/BF02105861}{DOI}.

\bibitem{Benois}
D. Benois, A generalization of Greenberg's $\mathcal L$-invariant,
\emph{Amer. J. Math.} \textbf{133} (2011), 1573--1632.
\href{https://doi.org/10.1353/ajm.2011.0043}{DOI}.

\bibitem{BertapelleEtAl}
A. Bertapelle, M. Candilera and V. Cristante,
Monodromy of logarithmic Barsotti--Tate groups attached to $1$-motives,
\emph{J. Reine Angew. Math.} \textbf{573} (2004), 211--234.
\href{https://arxiv.org/abs/math/0307059}{arXiv}.

\bibitem{BesserShalit}
A. Besser and E. de Shalit, $\mathcal L$-invariants of $p$-adically
uniformized varieties, \emph{Ann. Math. Qu\'ebec} \textbf{40} (2016), 29--54.
\href{https://doi.org/10.1007/s40316-015-0047-1}{DOI}.

\bibitem{BlochKato}
S. Bloch and K. Kato, $L$-functions and Tamagawa numbers of motives,
in \emph{The Grothendieck Festschrift}, Vol.~I, Progr. Math.,
vol.~86, Birkh\"auser, Boston, 1990, pp.~333--400.

\bibitem{BCDT}
C. Breuil, B. Conrad, F. Diamond and R. Taylor,
On the modularity of elliptic curves over $\Q$: wild $3$-adic exercises,
\emph{J. Amer. Math. Soc.} \textbf{14} (2001), 843--939.

\bibitem{Brumer}
A. Brumer, On the units of algebraic number fields,
\emph{Mathematika} \textbf{14} (1967), 121--124.

\bibitem{CM}
F. Calegari and B. Mazur, Nearly ordinary Galois deformations over
arbitrary number fields, \emph{J. Inst. Math. Jussieu}
\textbf{8} (2009), 99--177.
\href{https://doi.org/10.1017/S1474748008000327}{DOI}.

\bibitem{Cassou}
P. Cassou-Nogu\`es, Valeurs aux entiers n\'egatifs des fonctions
z\^eta et fonctions z\^eta $p$-adiques, \emph{Invent. Math.}
\textbf{51} (1979), 29--59.
\href{https://doi.org/10.1007/BF01389911}{DOI}.

\bibitem{CD}
P. Charollois and S. Dasgupta, Integral Eisenstein cocycles on
$\mathrm{GL}_n$, I: Sczech's cocycle and $p$-adic $L$-functions
of totally real fields, \emph{Camb. J. Math.} \textbf{2} (2014), 49--90.
\href{https://doi.org/10.4310/CJM.2014.v2.n1.a2}{DOI}.

\bibitem{CN}
M. Cimpoea\c{s} and F. Nicolae, Independence of Artin $L$-functions,
\emph{Forum Math.} \textbf{31} (2019), 529--534.
\href{https://arxiv.org/abs/1805.07990}{arXiv}.

\bibitem{Colmez}
P. Colmez, R\'esidu en $s=1$ des fonctions z\^eta $p$-adiques,
\emph{Invent. Math.} \textbf{91} (1988), 371--389.
\href{https://webusers.imj-prg.fr/~pierre.colmez/residu.pdf}{Author's copy}.

\bibitem{DHS}
S. Dasgupta, M. H. L. Honnor and M. Spie{\ss}, On the equality of
three formulas for Brumer--Stark units, \emph{J. Lond. Math. Soc.}
(2) \textbf{112} (2025), e70296.
\href{https://doi.org/10.1112/jlms.70296}{DOI}.

\bibitem{DKIntegral}
S. Dasgupta and M. Kakde, Brumer--Stark units and explicit class
field theory, \emph{Duke Math. J.} \textbf{173} (2024), 1477--1555.
\href{https://doi.org/10.1215/00127094-2023-0039}{DOI}.

\bibitem{DKV}
S. Dasgupta, M. Kakde and K. Ventullo, On the Gross--Stark conjecture,
\emph{Ann. of Math.} (2) \textbf{188} (2018), 833--870.
\href{https://annals.math.princeton.edu/2018/188-3/p03}{Journal}.

\bibitem{DS}
S. Dasgupta and M. Spiess,
On the characteristic polynomial of the Gross regulator matrix,
\emph{Trans. Amer. Math. Soc.} \textbf{371} (2019), 803--827.
\href{https://doi.org/10.1090/tran/7393}{DOI}.

\bibitem{DeligneIII}
P. Deligne, Th\'eorie de Hodge. III,
\emph{Publ. Math. Inst. Hautes \'Etudes Sci.} \textbf{44} (1974), 5--77.
\href{http://www.numdam.org/item/PMIHES_1974__44__5_0/}{Numdam}.

\bibitem{DeligneL}
P. Deligne, Valeurs de fonctions $L$ et p\'eriodes d'int\'egrales,
in \emph{Automorphic Forms, Representations and $L$-Functions},
Part~2, Proc. Sympos. Pure Math., vol.~33, Amer. Math. Soc.,
Providence, RI, 1979, pp.~313--346.

\bibitem{DR}
P. Deligne and K. A. Ribet, Values of abelian $L$-functions at
negative integers over totally real fields,
\emph{Invent. Math.} \textbf{59} (1980), 227--286.

\bibitem{FuchsPham}
C. Fuchs and D. H. Pham,
Some applications of the $p$-adic analytic subgroup theorem,
\emph{Glas. Mat. Ser. III} \textbf{51} (2016), 335--343.
\href{https://doi.org/10.3336/gm.51.2.04}{DOI}.

\bibitem{GreenbergArtin}
R. Greenberg, On $p$-adic Artin $L$-functions,
\emph{Nagoya Math. J.} \textbf{89} (1983), 77--87.
\href{https://doi.org/10.1017/S0027763000020250}{DOI}.

\bibitem{GS}
R. Greenberg and G. Stevens, $p$-adic $L$-functions and $p$-adic
periods of modular forms, \emph{Invent. Math.} \textbf{111}
(1993), 407--447.
\href{https://doi.org/10.1007/BF01231294}{DOI}.

\bibitem{Gross}
B. H. Gross, $p$-adic $L$-series at $s=0$,
\emph{J. Fac. Sci. Univ. Tokyo Sect. IA Math.}
\textbf{28} (1981), 979--994 (published 1982).

\bibitem{HK}
M. Hofer and S. Kleine, On two conjectures of Gross for large vanishing
orders, \emph{Monatsh. Math.} \textbf{197} (2022), 85--109.
\href{https://doi.org/10.1007/s00605-021-01555-3}{DOI}.

\bibitem{Honnor}
M. H. Honnor, Comparing two formulas for the Gross--Stark units,
\emph{J. Number Theory} \textbf{248} (2023), 368--419.
\href{https://doi.org/10.1016/j.jnt.2023.02.002}{DOI}.

\bibitem{HW}
A. Huber and G. W\"ustholz,
\emph{Transcendence and Linear Relations of $1$-Periods},
Cambridge Tracts in Mathematics, vol.~227,
Cambridge University Press, Cambridge, 2022.
\href{https://arxiv.org/abs/1805.10104}{arXiv}.

\bibitem{KMP}
J. Kaczorowski, G. Molteni and A. Perelli, Linear independence of
$L$-functions, \emph{Forum Math.} \textbf{18} (2006), 1--7.
\href{https://doi.org/10.1515/FORUM.2006.001}{DOI}.

\bibitem{Kobayashi}
S. Kobayashi, An elementary proof of the Mazur--Tate--Teitelbaum
conjecture for elliptic curves, \emph{Doc. Math.}, Extra Volume
Coates (2006), 567--575.
\href{https://doi.org/10.4171/DMS/4/16}{DOI}.

\bibitem{Lu}
F. L\"u, Independence of Dirichlet series,
\emph{Proc. Amer. Math. Soc.} \textbf{153} (2025), 3303--3313.
\href{https://doi.org/10.1090/proc/17276}{DOI}.

\bibitem{Maksoud}
A. Maksoud, On the rank of Leopoldt's and Gross's regulator maps,
\emph{Doc. Math.} \textbf{28} (2023), 1441--1471.
\href{https://ems.press/content/serial-article-files/31139}{Journal}.

\bibitem{MTT}
B. Mazur, J. Tate and J. Teitelbaum,
On $p$-adic analogues of the conjectures of Birch and Swinnerton-Dyer,
\emph{Invent. Math.} \textbf{84} (1986), 1--48.

\bibitem{Milne}
J. S. Milne, \emph{Arithmetic Duality Theorems}, 2nd ed.,
BookSurge, 2006.
\href{https://www.jmilne.org/math/Books/adt.html}{Author's edition}.

\bibitem{MohajerThesis}
M. Mohajer, \emph{Linear Relations of $p$-adic Periods of $1$-Motives},
doctoral thesis, University of Ottawa, 2024.
\href{https://doi.org/10.20381/ruor-30802}{Repository}.

\bibitem{MohajerSemistable}
M. Mohajer, Semistable $1$-Motives: Monodromy, Explicit Periods,
and Local Duality, preprint, 2026.
\href{https://arxiv.org/abs/2611.17685}{arXiv:2611.17685}.

\bibitem{MS}
M. Mohajer and A. Sebbar, $p$-adic Period Conjectures for $1$-Motives:
Integration and Linear Relations, preprint, 2025.
\href{https://arxiv.org/abs/2507.15020}{arXiv:2507.15020}.

\bibitem{Nekovar}
J. Nekov\'a\v r, Beilinson's conjectures, in \emph{Motives}
(Seattle, WA, 1991), Part~1, Proc. Sympos. Pure Math., vol.~55,
Amer. Math. Soc., Providence, RI, 1994, pp.~537--570.

\bibitem{Neukirch}
J. Neukirch, \emph{Algebraic Number Theory},
Grundlehren Math. Wiss., vol.~322, Springer, Berlin, 1999.

\bibitem{PerrinRiou}
B. Perrin-Riou, Fonctions $L$ $p$-adiques des repr\'esentations
$p$-adiques, \emph{Ast\'erisque} \textbf{229} (1995), 1--198.
\href{http://www.numdam.org/item/AST_1995__229__1_0/}{Numdam}.

\bibitem{Raynaud}
M. Raynaud, $1$-motifs et monodromie g\'eom\'etrique, in
\emph{P\'eriodes $p$-adiques}, \emph{Ast\'erisque} \textbf{223}
(1994), Expos\'e VII, 295--319.
\href{http://www.numdam.org/item/AST_1994__223__295_0/}{Numdam}.

\bibitem{Rivera}
C. Rivera Arredondo, Height pairings of $1$-motives,
\emph{Rend. Semin. Mat. Univ. Padova} \textbf{149} (2023), 151--181.
\href{https://www.numdam.org/item/RSMUP_2023__149__151_0/}{Numdam}.

\bibitem{Scholl}
A. J. Scholl, Remarks on special values of $L$-functions,
in \emph{$L$-Functions and Arithmetic} (Durham, 1989),
London Math. Soc. Lecture Note Ser., vol.~153, Cambridge University
Press, Cambridge, 1991, pp.~373--392.
\href{https://doi.org/10.1017/CBO9780511526053.016}{DOI}.

\bibitem{Spiess}
M. Spiess, Shintani cocycles and the order of vanishing of $p$-adic
Hecke $L$-series at $s=0$,
\emph{Math. Ann.} \textbf{359} (2014), 239--265.
\href{https://arxiv.org/abs/1203.6689}{arXiv}.

\bibitem{Stevens}
G. Stevens, \emph{Arithmetic on Modular Curves},
Progr. Math., vol.~20, Birkh\"auser, Boston, 1982.

\bibitem{TaylorWiles}
R. Taylor and A. Wiles, Ring-theoretic properties of certain Hecke
algebras, \emph{Ann. of Math.} (2) \textbf{141} (1995), 553--572.

\bibitem{Washington}
L. C. Washington, \emph{Introduction to Cyclotomic Fields}, 2nd ed.,
Graduate Texts in Mathematics, vol.~83, Springer, New York, 1997.
\href{https://doi.org/10.1007/978-1-4612-1934-7}{DOI}.

\bibitem{Wiles}
A. Wiles, Modular elliptic curves and Fermat's last theorem,
\emph{Ann. of Math.} (2) \textbf{141} (1995), 443--551.

\bibitem{Yoshikawa}
S. Yoshikawa, Independence of Hasse--Weil $L$-functions over rings
of Dirichlet series, preprint, 2026.
\href{https://arxiv.org/abs/2611.12541}{arXiv:2611.12541}.

\end{thebibliography}

\end{document}

